\documentclass[11pt]{article}
% Shared preamble for all solution files.
\usepackage[margin=1.1in]{geometry}
\usepackage{enumitem}
\usepackage{tikz}
\usepackage[hidelinks]{hyperref}
\usepackage{amsmath,amssymb,amsthm}
\newcommand{\fa}{\mathfrak a}
\newcommand{\fb}{\mathfrak b}
\newcommand{\fc}{\mathfrak c}
\newcommand{\fp}{\mathfrak p}
\newcommand{\fq}{\mathfrak q}
\newcommand{\fm}{\mathfrak m}
\newcommand{\fn}{\mathfrak n}
\newcommand{\fr}{\mathfrak r}
\newcommand{\Nil}{\mathfrak N}
\newcommand{\Jac}{\mathfrak R}
\newcommand{\rad}{r}
\newcommand{\Spec}{\operatorname{Spec}}
\newcommand{\Max}{\operatorname{Max}}
\newcommand{\Ann}{\operatorname{Ann}}
\newcommand{\Ker}{\operatorname{Ker}}
\newcommand{\Z}{\mathbf Z}
\newcommand{\Q}{\mathbf Q}
\newcommand{\R}{\mathbf R}
\newcommand{\C}{\mathbf C}
\newcommand{\ci}{\mathrm{c}}
\newcommand{\ex}{\mathrm{e}}

\setlength{\parindent}{0pt}
\setlength{\parskip}{0.55em}
% \exercise{label}{statement}  -- followed by the solution text
\newcommand{\exercise}[2]{\par\bigskip\noindent{\large\textbf{Exercise #1}}\par\smallskip
  \begin{quote}\small\textbf{Problem.}\ #2\end{quote}\par}
\newcommand{\pf}{\noindent\textbf{Solution.}\ }

\title{Liu, Chapter 1: Some topics in commutative algebra\\[2pt]\large Solutions to the exercises}
\author{}
\date{}
\begin{document}
\maketitle
\noindent\textbf{Conventions.} Rings are commutative with $1$. Results of the text are cited by their numbers in Liu (``Prop.~1.5'', ``Thm.~2.16'', \dots). The numbering and statements of the exercises follow the book (Exercises 1.1--1.7, 2.1--2.19, 3.1--3.15); each exercise is restated before its solution. A few sub-items that the book displays as nested lists are written as running prose.

Honest notes. Ex.~2.9 is false as printed for $M=0$ (zero is flat but its localizations are free of rank $0$), so it is proved for $M\neq0$. Ex.~3.6 is proved for a nonzero proper ideal $I$ (for $I=A$ the conclusion is trivial). Ex.~1.6 uses a directed index set, and Ex.~1.4(b) is read as ``finitely generated as an $A$-algebra''. In Ex.~3.3 the Taylor map on $C^\infty$ germs is not injective, and the solution says what is true. Some solutions of Section 1.3 depend on earlier ones (3.7 uses 3.3; 3.10 uses 3.8, 3.9; 3.12 uses 3.10).

\section*{1.1 Tensor products}

Throughout, $A$ is a commutative ring with unit, all tensor products without subscript are over $A$, and results of the text are cited by their numbers in Liu (``Prop.~1.5'', ``Cor.~1.13'', \dots).

\exercise{1.1}{Let $\{M_i\}_{i\in I}$, $\{N_j\}_{j\in J}$ be two families of modules over a ring $A$. Show that
\[ \Bigl(\bigoplus_{i\in I}M_i\Bigr)\otimes_A\Bigl(\bigoplus_{j\in J}N_j\Bigr)\simeq\bigoplus_{(i,j)\in I\times J}(M_i\otimes_A N_j). \]}
\pf
Write $M=\bigoplus_i M_i$ and $N=\bigoplus_j N_j$. We chain canonical isomorphisms from Prop.~1.5:
\begin{align*}
M\otimes N&\simeq\bigoplus_{i\in I}(M_i\otimes N) &&\text{by 1.5(d)},\\
M_i\otimes N&\simeq N\otimes M_i &&\text{by 1.5(b)},\\
N\otimes M_i&\simeq\bigoplus_{j\in J}(N_j\otimes M_i) &&\text{by 1.5(d)},\\
N_j\otimes M_i&\simeq M_i\otimes N_j &&\text{by 1.5(b)}.
\end{align*}
Each isomorphism in the second to fourth lines is applied inside a direct sum, componentwise (a direct sum of isomorphisms is an isomorphism). Composing, we get an isomorphism
\[ M\otimes N\;\xrightarrow{\ \sim\ }\;\bigoplus_{i\in I}\bigoplus_{j\in J}(M_i\otimes N_j)=\bigoplus_{(i,j)\in I\times J}(M_i\otimes N_j). \]
Following a pure tensor $x\otimes y$ with $x=(x_i)_i$, $y=(y_j)_j$ through the chain, it goes to the family whose $(i,j)$-component is $x_i\otimes y_j$. So the isomorphism is the canonical one, characterized by sending $\iota_i(x_i)\otimes\iota_j(y_j)$ to the element $x_i\otimes y_j$ placed in the $(i,j)$-summand. \hfill$\square$

\exercise{1.2}{Show that there exists a unique $A$-linear map
\[ f:\operatorname{Hom}_A(M,M')\otimes_A\operatorname{Hom}_A(N,N')\to\operatorname{Hom}_A(M\otimes_A N,M'\otimes_A N') \]
such that $f(u\otimes v)=u\otimes v$.}
\pf
On the right, $u\otimes v$ denotes the map of Def.~1.8, the unique $A$-linear map $M\otimes N\to M'\otimes N'$ with $x\otimes y\mapsto u(x)\otimes v(y)$.

\emph{Existence.} Define $g:\operatorname{Hom}_A(M,M')\times\operatorname{Hom}_A(N,N')\to\operatorname{Hom}_A(M\otimes N,M'\otimes N')$ by $g(u,v)=u\otimes v$. We check that $g$ is $A$-bilinear. Let $u,u'\in\operatorname{Hom}_A(M,M')$, $v\in\operatorname{Hom}_A(N,N')$, $a\in A$. For $x\in M$, $y\in N$,
\[ ((u+u')\otimes v)(x\otimes y)=(u(x)+u'(x))\otimes v(y)=(u\otimes v)(x\otimes y)+(u'\otimes v)(x\otimes y), \]
\[ ((au)\otimes v)(x\otimes y)=a\,u(x)\otimes v(y)=a\,(u\otimes v)(x\otimes y). \]
Both sides of each identity are $A$-linear maps on $M\otimes N$, and the elements $x\otimes y$ generate $M\otimes N$ (Rem.~1.3). Hence $(u+u')\otimes v=u\otimes v+u'\otimes v$ and $(au)\otimes v=a(u\otimes v)$. Linearity in $v$ is identical. By the universal property of the tensor product (Def.~1.1) there is a unique $A$-linear $f$ with $f(u\otimes v)=g(u,v)=u\otimes v$ for all $u,v$.

\emph{Uniqueness.} The elements $u\otimes v$ generate the source (Rem.~1.3), so an $A$-linear map is determined by its values on them. This is also part of the uniqueness in Def.~1.1. \hfill$\square$

\exercise{1.3}{Let $M,N$ be $A$-modules, and $i:M'\to M$, $j:N'\to N$ submodules of $M$ and $N$, respectively. Then there exists a canonical isomorphism
\[ (M/M')\otimes_A(N/N')\simeq (M\otimes_A N)/(\operatorname{Im}i_N+\operatorname{Im}j_M). \]
Show that $(\mathbb Z/n\mathbb Z)\otimes_{\mathbb Z}(\mathbb Z/m\mathbb Z)=\mathbb Z/l\mathbb Z$, where $l=\gcd(m,n)$.}
\pf
Here $i_N=i\otimes\operatorname{Id}_N:M'\otimes N\to M\otimes N$ and $j_M=j\otimes\operatorname{Id}_M:N'\otimes M\to N\otimes M\simeq M\otimes N$. Hence $\operatorname{Im}i_N$ is the submodule of $M\otimes N$ generated by the $x'\otimes y$ ($x'\in M'$, $y\in N$), and $\operatorname{Im}j_M$ is generated by the $x\otimes y'$ ($x\in M$, $y'\in N'$). Put $K=\operatorname{Im}i_N+\operatorname{Im}j_M$.

\emph{Step 1.} Let $\pi:M\to M/M'$ be the projection and $\pi_N=\pi\otimes\operatorname{Id}_N:M\otimes N\to(M/M')\otimes N$. By Cor.~1.13 (applied to $M'\to M$ and the module $N$), $\pi_N$ is surjective with kernel $\operatorname{Im}i_N$.

\emph{Step 2.} Let $q:N\to N/N'$ be the projection. The sequence $N'\xrightarrow{j}N\xrightarrow{q}N/N'\to0$ is exact, so by Prop.~1.12 (tensoring with $M/M'$, and using commutativity)
\[ (M/M')\otimes N'\xrightarrow{\ \operatorname{Id}\otimes j\ }(M/M')\otimes N\xrightarrow{\ \operatorname{Id}\otimes q\ }(M/M')\otimes(N/N')\to0 \]
is exact. Thus $\operatorname{Id}\otimes q$ is surjective with kernel $W=\operatorname{Im}(\operatorname{Id}\otimes j)$. This kernel is generated by the elements $\bar x\otimes y'$ with $x\in M$, $y'\in N'$, which are the images under $\pi_N$ of the generators $x\otimes y'$ of $\operatorname{Im}j_M$. Hence $W=\pi_N(\operatorname{Im}j_M)$.

\emph{Step 3.} Consider the surjection $\Theta=(\operatorname{Id}\otimes q)\circ\pi_N:M\otimes N\to(M/M')\otimes(N/N')$, so that $\Theta(x\otimes y)=\bar x\otimes\bar y$. Its kernel is $\pi_N^{-1}(W)=\pi_N^{-1}(\pi_N(\operatorname{Im}j_M))=\operatorname{Im}j_M+\ker\pi_N=\operatorname{Im}j_M+\operatorname{Im}i_N=K$. Therefore $\Theta$ induces an isomorphism
\[ (M\otimes N)/K\xrightarrow{\ \sim\ }(M/M')\otimes(N/N'),\qquad \overline{x\otimes y}\mapsto\bar x\otimes\bar y, \]
which is canonical.

\emph{The case of $\mathbb Z$.} Take $A=\mathbb Z$. By Cor.~1.13 (with $I=n\mathbb Z$ and $M=\mathbb Z/m\mathbb Z$, and Prop.~1.5(b)),
\[ (\mathbb Z/n\mathbb Z)\otimes(\mathbb Z/m\mathbb Z)\simeq(\mathbb Z/m\mathbb Z)\big/n(\mathbb Z/m\mathbb Z). \]
The subgroup $n(\mathbb Z/m\mathbb Z)$ of $\mathbb Z/m\mathbb Z$ is $(n\mathbb Z+m\mathbb Z)/m\mathbb Z=l\mathbb Z/m\mathbb Z$, because $n\mathbb Z+m\mathbb Z=l\mathbb Z$. So the quotient is $(\mathbb Z/m\mathbb Z)/(l\mathbb Z/m\mathbb Z)\simeq\mathbb Z/l\mathbb Z$. (When $l=1$ this recovers Example~1.4.) \hfill$\square$

\exercise{1.4}{Let $M,N$ be $A$-modules, and let $B,C$ be $A$-algebras.
\begin{enumerate}[label=(\alph*),leftmargin=*]
\item If $M$ and $N$ are finitely generated over $A$, then so is $M\otimes_A N$.
\item If $B$ and $C$ are finitely generated over $A$, then so is $B\otimes_A C$.
\item Taking $A=\mathbb Z$, $M=B=\mathbb Z/2\mathbb Z$, and $N=C=\mathbb Q$, show that the converse of (a) and (b) is false.
\end{enumerate}}
\pf
In (b), ``finitely generated over $A$'' means finitely generated as an $A$-algebra. (If $B,C$ are finitely generated as $A$-modules, then so is $B\otimes_A C$ by (a); this is a stronger hypothesis and a stronger conclusion, which gives the module version of (b).)

(a) Let $x_1,\dots,x_r$ generate $M$ and $y_1,\dots,y_s$ generate $N$. By Rem.~1.3, $M\otimes N$ is generated by the elements $x\otimes y$. Write $x=\sum_k a_kx_k$ and $y=\sum_l b_ly_l$. Then $x\otimes y=\sum_{k,l}a_kb_l\,(x_k\otimes y_l)$. So the $rs$ elements $x_k\otimes y_l$ generate $M\otimes N$.

(b) Let $B=A[b_1,\dots,b_r]$ and $C=A[c_1,\dots,c_s]$ be generated as algebras. Let $D\subseteq B\otimes C$ be the $A$-subalgebra generated by the $b_k\otimes1=p_1(b_k)$ and $1\otimes c_l=p_2(c_l)$. Since $p_1,p_2$ are $A$-algebra homomorphisms, $D$ contains $p_1(B)$ and $p_2(C)$: for a polynomial $P$ with coefficients in $A$, $p_1(P(b_1,\dots,b_r))=P(b_1\otimes1,\dots,b_r\otimes1)\in D$, and similarly for $p_2$. Hence for $b\in B$, $c\in C$ we get $b\otimes c=(b\otimes1)(1\otimes c)=p_1(b)p_2(c)\in D$. Since the elements $b\otimes c$ generate $B\otimes C$ additively (Rem.~1.3), $D=B\otimes C$. So $B\otimes_A C$ is generated by the $r+s$ elements $b_k\otimes 1$, $1\otimes c_l$.

(c) Let $A=\mathbb Z$. For $x\in\mathbb Z/2\mathbb Z$ and $y\in\mathbb Q$ we have $x\otimes y=x\otimes 2(y/2)=(2x)\otimes(y/2)=0$, since $2x=0$. As such elements generate (Rem.~1.3), $M\otimes N=(\mathbb Z/2\mathbb Z)\otimes_{\mathbb Z}\mathbb Q=0$. The same computation gives $B\otimes_{\mathbb Z}C=0$. The zero module is finitely generated over $\mathbb Z$, and the zero ring is generated as a $\mathbb Z$-algebra by $\emptyset$ (or by $0$). Also $M=B=\mathbb Z/2\mathbb Z$ is finitely generated as a module (by $1$) and as an algebra. But $\mathbb Q$ is neither:
\begin{itemize}
\item \emph{$\mathbb Q$ is not a finitely generated $\mathbb Z$-module.} If $q_1,\dots,q_r\in\mathbb Q$ generate $\mathbb Q$, choose a common denominator $d\ge1$ with $dq_k\in\mathbb Z$ for all $k$. Then every element of $\mathbb Q$ lies in $\frac1d\mathbb Z$. But $\frac1{2d}\notin\frac1d\mathbb Z$, a contradiction.
\item \emph{$\mathbb Q$ is not a finitely generated $\mathbb Z$-algebra.} Suppose $\mathbb Q=\mathbb Z[a_1/b_1,\dots,a_r/b_r]$ with $a_k,b_k\in\mathbb Z$, $b_k\neq0$. Put $b=b_1\cdots b_r$. Then $\mathbb Q\subseteq\mathbb Z[1/b]=\{m/b^t:m\in\mathbb Z,\ t\ge0\}$. Take a prime $p\nmid b$. If $1/p=m/b^t$, then $b^t=pm$, so $p\mid b^t$ and hence $p\mid b$, a contradiction. So $1/p\notin\mathbb Z[1/b]$.
\end{itemize}
Therefore $M\otimes N$ and $B\otimes C$ are finitely generated while $N$ and $C$ are not, so the converses of (a) and (b) fail. \hfill$\square$

\exercise{1.5}{Let $\rho:A\to B$ be a ring homomorphism, $M$ an $A$-module, and $N$ a $B$-module. Show that there exists a canonical isomorphism of $A$-modules
\[ \operatorname{Hom}_A(M,\rho_*N)\simeq\operatorname{Hom}_B(\rho^*M,N). \]}
\pf
Recall $\rho_*N$ is $N$ with $a\cdot y=\rho(a)y$, and $\rho^*M=M\otimes_AB$ with $B$-action $b'(x\otimes b)=x\otimes b'b$. The right side is a $B$-module, hence an $A$-module through $\rho$; the left side is an $A$-module in the usual way.

\emph{The map $\Phi$.} Let $g\in\operatorname{Hom}_A(M,\rho_*N)$. The map $M\times B\to N$, $(x,b)\mapsto b\,g(x)$, is $A$-bilinear. It is biadditive. It is compatible with scalars: for $a\in A$, $(ax,b)\mapsto b\,g(ax)=b\rho(a)g(x)=(\rho(a)b)g(x)$ and $(x,\rho(a)b)\mapsto(\rho(a)b)g(x)$, and $(x,ab)\mapsto\rho(a)b\,g(x)$ also (here $a\cdot b=\rho(a)b$ in $B$). So it induces an $A$-linear map $\tilde g:M\otimes_AB\to N$ with $\tilde g(x\otimes b)=b\,g(x)$. It is $B$-linear: $\tilde g(b'(x\otimes b))=\tilde g(x\otimes b'b)=b'b\,g(x)=b'\tilde g(x\otimes b)$, and the identity on generators suffices. Put $\Phi(g)=\tilde g$.

\emph{The map $\Psi$.} For $h\in\operatorname{Hom}_B(\rho^*M,N)$ put $\Psi(h)(x)=h(x\otimes1)$. It is $A$-linear into $\rho_*N$: additivity is clear and
\[ h(ax\otimes1)=h(x\otimes\rho(a))=h(\rho(a)(x\otimes1))=\rho(a)h(x\otimes1)=a\cdot h(x\otimes 1). \]

\emph{Inverse.} $\Psi\Phi(g)(x)=\tilde g(x\otimes1)=g(x)$. And $\Phi\Psi(h)(x\otimes b)=b\,h(x\otimes1)=h(b(x\otimes1))=h(x\otimes b)$; since the $x\otimes b$ generate $M\otimes_AB$, $\Phi\Psi(h)=h$. So $\Phi$ is bijective with inverse $\Psi$.

\emph{$A$-linearity.} $\Phi$ is clearly additive. For $a\in A$, $\Phi(ag)(x\otimes b)=b\,\rho(a)g(x)=\rho(a)\,\Phi(g)(x\otimes b)=(a\Phi(g))(x\otimes b)$. Hence $\Phi$ is an isomorphism of $A$-modules. It involves no choices, so it is canonical. \hfill$\square$

\exercise{1.6}{Let $(N_i)_{i\in I}$ be a direct system of $A$-modules. Then for any $A$-module $M$, there exists a canonical isomorphism $\varinjlim(N_i\otimes_AM)\simeq(\varinjlim N_i)\otimes_AM$.}
\pf
As usual, a direct system is indexed by a \emph{directed} set $I$ (for $i,j\in I$ there is $k\ge i,j$), with transition maps $\varphi_{ik}:N_i\to N_k$ for $i\le k$. (For a non-directed index set the statement is still true, but the argument below, which uses the elementwise description of the limit, needs directedness.) We use the standard facts about $N=\varinjlim N_i$ with canonical maps $\lambda_i:N_i\to N$: every element of $N$ is of the form $\lambda_i(n_i)$; and $\lambda_i(n_i)=\lambda_j(n_j)$ if and only if $\varphi_{ik}(n_i)=\varphi_{jk}(n_j)$ for some $k\ge i,j$. Let $P=\varinjlim(N_i\otimes M)$, for the system with transition maps $\varphi_{ik}\otimes\operatorname{Id}_M$, and canonical maps $\mu_i:N_i\otimes M\to P$. Every element of $P$ is of the form $\mu_i(z)$ with $z\in N_i\otimes M$.

\emph{The map $\alpha:P\to N\otimes M$.} The maps $\lambda_i\otimes\operatorname{Id}_M:N_i\otimes M\to N\otimes M$ are compatible with the transition maps, because $\lambda_k\circ\varphi_{ik}=\lambda_i$. By the universal property of the direct limit they induce $\alpha$ with $\alpha\mu_i(n_i\otimes m)=\lambda_i(n_i)\otimes m$.

\emph{The map $\beta:N\otimes M\to P$.} Define $b:N\times M\to P$ by $b(n,m)=\mu_i(n_i\otimes m)$ where $n=\lambda_i(n_i)$. This is well defined. If also $n=\lambda_j(n_j)$, pick $k\ge i,j$ with $\varphi_{ik}(n_i)=\varphi_{jk}(n_j)$. Then
\[ \mu_i(n_i\otimes m)=\mu_k(\varphi_{ik}(n_i)\otimes m)=\mu_k(\varphi_{jk}(n_j)\otimes m)=\mu_j(n_j\otimes m), \]
using $\mu_k\circ(\varphi_{ik}\otimes\operatorname{Id})=\mu_i$. It is bilinear. Given $n,n'\in N$, write them as $\lambda_i(n_i),\lambda_i(n'_i)$ with a common index $i$ (possible by directedness, after pushing forward to a common $k$). Then $n+n'=\lambda_i(n_i+n'_i)$, so $b(n+n',m)=\mu_i((n_i+n_i')\otimes m)=b(n,m)+b(n',m)$. Additivity in $m$ and compatibility with scalars $a\in A$ follow in the same way from $a n=\lambda_i(an_i)$ and the bilinearity of $\otimes$. So $b$ induces an $A$-linear $\beta$ with $\beta(n\otimes m)=\mu_i(n_i\otimes m)$.

\emph{Inverse isomorphisms.} For $n_i\in N_i$, $m\in M$ we have $\beta\alpha\mu_i(n_i\otimes m)=\beta(\lambda_i(n_i)\otimes m)=\mu_i(n_i\otimes m)$. The elements $\mu_i(n_i\otimes m)$ generate $P$ (every element of $P$ is some $\mu_i(z)$ with $z$ a sum of pure tensors), so $\beta\alpha=\operatorname{Id}_P$. Also $\alpha\beta(n\otimes m)=\alpha\mu_i(n_i\otimes m)=\lambda_i(n_i)\otimes m=n\otimes m$, and the $n\otimes m$ generate $N\otimes M$, so $\alpha\beta=\operatorname{Id}$. Hence $\alpha$ is an isomorphism, and it is canonical (it sends $\mu_i(n_i\otimes m)$ to $\lambda_i(n_i)\otimes m$).

\emph{Remark on the hint.} Equivalently, $P$ together with the bilinear map $b$ satisfies the universal property of $N\otimes M$, and the computation above is exactly the verification. \hfill$\square$

\exercise{1.7}{Let $B$ be an $A$-algebra, and let $M,N$ be $B$-modules. Show that there exists a canonical surjective homomorphism $M\otimes_AN\to M\otimes_BN$.}
\pf
Let $\rho:A\to B$ be the structure map; $M,N$ are $A$-modules via $\rho$, so $ax=\rho(a)x$. Consider the map $\theta:M\times N\to M\otimes_BN$, $\theta(x,y)=x\otimes_By$. It is biadditive, and for $a\in A$,
\[ \theta(ax,y)=\rho(a)x\otimes_By=\rho(a)(x\otimes_By)=x\otimes_B\rho(a)y=\theta(x,ay)=a\,\theta(x,y). \]
So $\theta$ is $A$-bilinear, and by Def.~1.1 there is a unique $A$-linear map $\Theta:M\otimes_AN\to M\otimes_BN$ with $\Theta(x\otimes_Ay)=x\otimes_By$. It is canonical, as it is characterized by this formula (it is even $B$-linear). The $B$-module structure on $M\otimes_AN$ is given by $b(x\otimes_Ay)=bx\otimes_Ay$, and $\Theta$ is $B$-linear: $\Theta(b(x\otimes_Ay))=bx\otimes_By=b\,\Theta(x\otimes_Ay)$. Hence the image of $\Theta$ is a $B$-submodule of $M\otimes_BN$ containing every $x\otimes_By$. A $B$-submodule containing a $B$-generating set is everything, so $\Theta$ is surjective. \hfill$\square$

\section*{1.2 Flatness}

Throughout, rings are commutative with $1$. We use the following elementary fact, which is not stated in the text.

\textbf{(D)} \emph{A direct summand of a flat module is flat.} Indeed, let $M=M_1\oplus M_2$ be flat and let $u:N'\to N$ be injective. By Prop.~1.5 the map $u\otimes\operatorname{id}_M:N'\otimes M\to N\otimes M$ is identified with the direct sum of the maps $N'\otimes M_i\to N\otimes M_i$; a direct sum of maps is injective only if each summand is injective, so $M_1$ and $M_2$ are flat.

\exercise{2.1}{Let $M$ be an $A$-module, $\Ann(M)=\{a\in A\mid aM=0\}$ its annihilator, and $I\subseteq\Ann(M)$ an ideal.
\begin{enumerate}[label=(\alph*),leftmargin=*]
\item Show that $M$ is endowed, in a natural way, with the structure of an $A/I$-module, and that $M\simeq M\otimes_A A/I$.
\item Let $N$ be another $A$-module such that $I\subseteq\Ann(N)$. Show that the canonical homomorphism $M\otimes_A N\to M\otimes_{A/I}N$ is an isomorphism.
\end{enumerate}}
\pf
(a) Define $(a+I)\cdot m=am$. This is well defined: if $a-a'\in I$ then $(a-a')m=0$. The module axioms for $A/I$ follow from those for $A$. By Cor.~1.13, $M\otimes_A A/I\simeq M/IM$, and $IM=0$ because $I\subseteq \Ann(M)$. So $M\otimes_AA/I\simeq M$ (explicitly, $m\otimes \bar a\mapsto am$).

(b) By Exercise~1.7 there is a canonical surjection $\alpha:M\otimes_AN\to M\otimes_{A/I}N$, $m\otimes n\mapsto m\otimes n$. Conversely, the map $M\times N\to M\otimes_AN$, $(m,n)\mapsto m\otimes n$, is $A$-bilinear, and it is $A/I$-bilinear for the structures of (a): $(\bar a m)\otimes n=am\otimes n=m\otimes an=m\otimes \bar a n$. Hence it induces an additive map $\beta:M\otimes_{A/I}N\to M\otimes_AN$ with $\beta(m\otimes n)=m\otimes n$. Now $\alpha\beta$ and $\beta\alpha$ are the identity on the generators $m\otimes n$ of the respective tensor products, hence are the identity. So $\alpha$ is an isomorphism. \hfill$\square$

\exercise{2.2}{Let $\rho:A\to B$ be a ring homomorphism, $S$ a multiplicative subset of $A$, and $T=\rho(S)$. Show that $T$ is a multiplicative subset of $B$, and that $T^{-1}B\simeq B\otimes_AS^{-1}A$ as $A$-algebras.}
\pf
Since $\rho(1)=1$ and $\rho(s)\rho(s')=\rho(ss')$, $T$ is multiplicative. Regard $B$ as an $A$-module. By Lemma~2.10 there is an isomorphism of $S^{-1}A$-modules $B\otimes_AS^{-1}A\simeq S^{-1}B$, $b\otimes a/s\mapsto \rho(a)b/s$, where $S^{-1}B$ is the localization of the $A$-module $B$: elements $b/s$ with $b/s=b'/s'$ iff $s''(s'b-sb')=0$ for some $s''\in S$, i.e.\ $\rho(s'')(\rho(s')b-\rho(s)b')=0$ in $B$. This is exactly the equivalence relation defining $T^{-1}B$ for $b/\rho(s)$ and $b'/\rho(s')$, with $\rho(s'')\in T$ (every element of $T$ is of the form $\rho(s'')$). Hence $b/s\mapsto b/\rho(s)$ is a well-defined bijection $S^{-1}B\to T^{-1}B$, compatible with addition (both have $b/s+b'/s'=(s'b+sb')/(ss')$). Composing, we get a bijection
\[ \theta: B\otimes_AS^{-1}A\to T^{-1}B,\qquad b\otimes \tfrac as\mapsto \tfrac{\rho(a)b}{\rho(s)}. \]
It is a ring homomorphism: the ring structure on $B\otimes_AS^{-1}A$ is given by $(b\otimes x)(b'\otimes x')=bb'\otimes xx'$, and $\theta$ is multiplicative on such pure tensors because $\frac{\rho(a)b}{\rho(s)}\cdot\frac{\rho(a')b'}{\rho(s')}=\frac{\rho(aa')bb'}{\rho(ss')}$; $\theta(1)=1$. It is compatible with the maps from $A$: $a\mapsto 1\otimes a/1\mapsto \rho(a)/1$, and $a\mapsto \rho(a)\otimes1$ also maps to $\rho(a)/1$ (the two $A$-algebra structures on $B\otimes_AS^{-1}A$ through $B$ and through $A$ coincide, as $\rho(a)\otimes 1=a\cdot(1\otimes1)=1\otimes a/1$). So $\theta$ is an isomorphism of $A$-algebras (indeed of $B$-algebras and of $S^{-1}A$-algebras). \hfill$\square$

\exercise{2.3}{Show that Nakayama's lemma is false for modules $M$ that are not finitely generated.}
\pf
Let $A=\Z_{(p)}=\{a/b\in\Q\mid p\nmid b\}$, a local ring with maximal ideal $\fm=pA=\operatorname{Rad}(A)$, and let $M=\Q$, regarded as an $A$-module. Then $\fm M=pM=M$, since $x=p\cdot(x/p)$ for all $x\in\Q$; but $M\neq0$. So the statement of Lemma~2.7 (with $I=\operatorname{Rad}(A)$) fails for $M$. (Note that $M$ is not finitely generated, as any finitely generated $A$-submodule of $\Q$ is of the form $p^{-k}A$.) \hfill$\square$

\exercise{2.4}{Let $I$ be a finitely generated ideal of $A$. Show that the following properties are equivalent: (i) $A/I$ is flat over $A$; (ii) $I=I^2$; (iii) there exists $e\in A$ such that $e^2=e$ and $I=eA$.}
\pf
(i)$\Rightarrow$(ii). Let $M=A/I$ be flat. By Thm.~2.4 the canonical map $I\otimes_AM\to IM$ is an isomorphism. Now $IM=I(A/I)=0$, while $I\otimes_AA/I\simeq I/I^2$ by Cor.~1.13. Hence $I/I^2=0$, i.e.\ $I=I^2$.

(ii)$\Rightarrow$(iii). Let $I=(x_1,\dots,x_n)$. Since $I=I^2$, we can write $x_i=\sum_j a_{ij}x_j$ with $a_{ij}\in I$. Let $\Delta=\det(\delta_{ij}-a_{ij})$. Multiplying the system $\sum_j(\delta_{ij}-a_{ij})x_j=0$ by the adjugate matrix gives $\Delta x_j=0$ for all $j$, so $\Delta I=0$. Expanding the determinant, $\Delta=1-e$ with $e\in I$. Then for every $x\in I$, $(1-e)x=0$, i.e.\ $x=ex\in eA$. Thus $I\subseteq eA$, and $eA\subseteq I$ since $e\in I$. Finally $e\in I$ gives $e=ee$, so $e^2=e$.

(iii)$\Rightarrow$(i). Let $e^2=e$, $I=eA$. Then $A=eA\oplus(1-e)A$ (any $a=ea+(1-e)a$, and if $ea=(1-e)b$ then multiplying by $e$ gives $ea=0$). The composite $(1-e)A\to A\to A/I$ is surjective (since $a\equiv (1-e)a$ mod $I$) and injective (if $(1-e)a=eb$, multiplying by $1-e$ gives $(1-e)a=0$). So $A/I\simeq (1-e)A$ is a direct summand of the free, hence flat (Prop.~2.2(a)), module $A$, and is flat by (D). \hfill$\square$

\exercise{2.5}{Let $A$ be an integral domain. Show that every $A$-module is flat if and only if $A$ is a field.}
\pf
If $A$ is a field, every $A$-module is a vector space, hence free, hence flat (Prop.~2.2(a)).

Conversely, suppose every $A$-module is flat. Let $a\in A$ be nonzero and $I=aA$. Then $A/I$ is flat and $I$ is finitely generated, so by Exercise~2.4, $I=I^2$. Thus $a=a^2b$ for some $b\in A$, i.e.\ $a(1-ab)=0$. As $A$ is a domain and $a\neq0$, $ab=1$. So every nonzero element is a unit, and $A$ (which is a domain, so $1\ne 0$) is a field. \hfill$\square$

\exercise{2.6}{Let $B$ be a flat $A$-algebra, with structure map $\rho:A\to B$.
\begin{enumerate}[label=(\alph*),leftmargin=*]
\item Show that for any finite family $\{I_\lambda\}_{\lambda\in\Lambda}$ of ideals of $A$, we have $\bigcap_{\lambda\in\Lambda}(I_\lambda B)=\bigl(\bigcap_{\lambda\in\Lambda}I_\lambda\bigr)B$.
\item Suppose $B$ is faithfully flat. Show that $IB\cap A=I$ for any ideal $I$ of $A$ (here $IB\cap A$ means $\rho^{-1}(IB)$).
\item Let $k$ be a field, $A=k[t,s]$ and $C=A[z]/(tz-s)$. By considering the ideals $tA$ and $sA$, show that $C$ is not flat over $A$.
\end{enumerate}}
\pf
(a) The inclusion $\supseteq$ is clear. For $\subseteq$ we may assume $\Lambda\neq\emptyset$ (for $\Lambda=\emptyset$ both sides are $B$) and, by induction on $|\Lambda|$, it suffices to treat two ideals $I,J$: indeed $\bigcap_{\lambda}I_\lambda B=I_1B\cap\bigl(\bigcap_{\lambda\ne1}I_\lambda B\bigr)=I_1B\cap\bigl(\bigcap_{\lambda\neq1}I_\lambda\bigr)B=(I_1\cap\bigcap_{\lambda\ne1}I_\lambda)B$ by the case of two ideals.

For two ideals, consider the exact sequence of $A$-modules
\[ 0\to A/(I\cap J)\xrightarrow{\ \ x\mapsto(x,x)\ \ }A/I\oplus A/J. \]
Tensoring with the flat $A$-module $B$ and using $A/\fa\otimes_AB\simeq B/\fa B$ (Cor.~1.13) and Prop.~1.5, we get an exact sequence
\[ 0\to B/(I\cap J)B\to B/IB\oplus B/JB,\qquad \bar b\mapsto(\bar b,\bar b). \]
If $b\in IB\cap JB$, then $\bar b$ maps to $0$, so $b\in(I\cap J)B$. Thus $IB\cap JB\subseteq(I\cap J)B$.

(b) Let $J=\rho^{-1}(IB)$, an ideal of $A$ with $I\subseteq J$ and $\rho(J)\subseteq IB$, hence $JB\subseteq IB$; also $IB\subseteq JB$. So $IB=JB$. Let $f:A/I\to A/J$ be the canonical surjection. Under the identifications $A/\fa\otimes_AB=B/\fa B$ (Cor.~1.13), $f\otimes\operatorname{id}_B$ is the canonical map $B/IB\to B/JB$, which is an isomorphism because $IB=JB$. Since $B$ is faithfully flat, Prop.~2.17(iii) shows that $f$ is an isomorphism, i.e.\ $J=I$. So $IB\cap A=I$.

(c) Since $s=tz$ in $C$, we have $C=k[t,s,z]/(tz-s)\simeq k[t,z]$ with $t\mapsto t$, $s\mapsto tz$ (eliminate the variable $s$). Suppose $C$ is flat over $A$. By (a), $(tA\cap sA)C=tC\cap sC$. Now $tA\cap sA=tsA$ ($A$ is a UFD and $t,s$ are non-associate primes), so the left side is $t\cdot tz\,C=t^2zC$. The right side is $tC\cap tzC=tzC$ (as $tzC\subseteq tC$). But $tz\notin t^2zC$: otherwise $tz=t^2zg$ for some $g\in k[t,z]$, so $1=tg$ in the domain $k[t,z]$, impossible. This contradiction shows that $C$ is not flat over $A$. \hfill$\square$

\exercise{2.7}{Give an example of a finitely generated flat module that is not free (over a suitable ring $A$).}
\pf
Let $k$ be a field, $A=k\times k$, $e=(1,0)$, and $M=eA=k\times 0$. Then $e^2=e$, $A=eA\oplus(1-e)A$ (as in Exercise~2.4), so $M$ is a direct summand of the free module $A$ and hence flat by Prop.~2.2(a) and (D). It is finitely generated (by $e$). It is not free: it is nonzero, so a free structure would have rank $\ge1$, and then $M$ would contain a copy of $A$, forcing $\Ann(M)=0$. But $(0,1)\cdot M=0$ with $(0,1)\ne0$. Hence $M$ is flat, finitely generated, and not free. (Note that $A$ is not local, in accordance with Thm.~2.16.) \hfill$\square$

\exercise{2.8}{Let $A$ be a Noetherian ring, $M$ a finitely generated $A$-module, $N$ an $A$-module, and $B$ a flat $A$-algebra. Consider the canonical homomorphism
\[ \rho:\Hom_A(M,N)\otimes_AB\to\Hom_B(M\otimes_AB,N\otimes_AB). \]
\begin{enumerate}[label=(\alph*),leftmargin=*]
\item Show that $\rho$ is an isomorphism if $M$ is free of finite rank.
\item Let $0\to K\to L\to M\to0$ be an exact sequence of $A$-modules. Show that we have a canonical exact sequence $0\to\Hom_A(M,N)\to\Hom_A(L,N)\to\Hom_A(K,N)$.
\item By taking for $L$ a free module of finite rank, show that $\rho$ is injective. By applying the injectivity to $K$, show that $\rho$ is an isomorphism.
\end{enumerate}}
\pf
Here $\rho(\varphi\otimes b)$ is the $B$-linear map $m\otimes b'\mapsto\varphi(m)\otimes bb'$. We write $\rho_M$ to indicate the dependence on $M$ ($N$, $B$ being fixed); $\rho_M$ is defined for every $A$-module $M$, and it is natural in $M$: if $u:M_1\to M_2$ is $A$-linear, the square formed by $\rho_{M_2},\rho_{M_1}$ and the maps $u^*\otimes\operatorname{id}_B$, $(u\otimes\operatorname{id}_B)^*$ commutes, as one checks on the generators $\varphi\otimes b$ of $\Hom_A(M_2,N)\otimes B$: both ways give $m\otimes b'\mapsto\varphi(u(m))\otimes bb'$.

(a) Let $M=A^n$ with basis $e_1,\dots,e_n$. Then $\Hom_A(M,N)\simeq N^n$ via $\varphi\mapsto(\varphi(e_i))_i$; $M\otimes B\simeq B^n$ has the $B$-basis $e_i\otimes1$, so $\Hom_B(M\otimes B,N\otimes B)\simeq(N\otimes B)^n$. Under these identifications $\rho$ becomes $N^n\otimes_AB\to(N\otimes_AB)^n$, $(n_i)_i\otimes b\mapsto(n_i\otimes b)_i$, which is the canonical isomorphism of Prop.~1.5 (tensor product commutes with finite direct sums). Hence $\rho$ is an isomorphism.

(b) Let $j:K\to L$ and $p:L\to M$. The map $p^*:\Hom(M,N)\to\Hom(L,N)$, $\psi\mapsto\psi p$, is injective because $p$ is surjective. Next $j^*p^*\psi=\psi pj=0$. Conversely, if $\varphi\in\Hom(L,N)$ satisfies $\varphi j=0$, then $\varphi$ vanishes on $K=\Ker p$ and so factors as $\varphi=\psi p$ with $\psi\in\Hom(M,N)$ (since $p$ induces $L/K\simeq M$). So $\Ker j^*=\Imm p^*$.

(c) Since $M$ is finitely generated, there is an exact sequence $0\to K\to L\to M\to0$ with $L=A^n$ free of finite rank; $K$ is finitely generated since $A$ is Noetherian. By (b) and flatness of $B$ (which preserves exactness) we get an exact row
\[ 0\to\Hom_A(M,N)\otimes B\to\Hom_A(L,N)\otimes B\to\Hom_A(K,N)\otimes B. \]
Also $0\to K\otimes B\to L\otimes B\to M\otimes B\to0$ is exact because $B$ is flat, so (b) over $B$ applied to the $B$-module $N\otimes B$ gives an exact row
\[ 0\to\Hom_B(M_B,N_B)\to\Hom_B(L_B,N_B)\to\Hom_B(K_B,N_B), \]
where $X_B=X\otimes_AB$. The maps $\rho_M,\rho_L,\rho_K$ form a commutative diagram between these rows (naturality), and $\rho_L$ is an isomorphism by (a).

\emph{Injectivity of $\rho_M$.} Let $x\in\Hom_A(M,N)\otimes B$ with $\rho_M(x)=0$ and let $x_L$ be its image in $\Hom_A(L,N)\otimes B$. Then $\rho_L(x_L)$ is the image of $\rho_M(x)=0$, so $x_L=0$ as $\rho_L$ is injective; by exactness of the first row (injectivity of the first arrow), $x=0$. This argument applies to every finitely generated $A$-module $M$ (with such a presentation); in particular to $K$, so $\rho_K$ is injective.

\emph{Surjectivity of $\rho_M$.} Let $y\in\Hom_B(M_B,N_B)$ and let $y_L$ be its image in $\Hom_B(L_B,N_B)$. Write $y_L=\rho_L(z)$ with $z\in\Hom_A(L,N)\otimes B$. The image $z_K$ of $z$ in $\Hom_A(K,N)\otimes B$ satisfies $\rho_K(z_K)=$ image of $y_L$ in $\Hom_B(K_B,N_B)$, which is $0$ by exactness of the second row. As $\rho_K$ is injective, $z_K=0$, so by exactness of the first row $z$ is the image of some $x\in\Hom_A(M,N)\otimes B$. Then $\rho_M(x)$ and $y$ have the same image $\rho_L(z)=y_L$ in $\Hom_B(L_B,N_B)$, and the map $\Hom_B(M_B,N_B)\to\Hom_B(L_B,N_B)$ is injective, so $\rho_M(x)=y$. Thus $\rho$ is an isomorphism. \hfill$\square$

\exercise{2.9}{Let $A$ be an integral domain, $K$ its field of fractions, and $M$ a finitely generated sub-$A$-module of $K$. Show that $M$ is flat if and only if it is locally free of rank $1$ (i.e.\ $M_\fp$ is free of rank $1$ over $A_\fp$ for every prime ideal $\fp$ of $A$).}
\pf
\emph{Correction.} For $M=0$ the statement is false as printed: $0$ is flat, but $M_\fp=0$ is free of rank $0$, not $1$. We prove the statement for $M\neq0$ (for $M=0$ the correct statement is that $M_\fp$ is free of rank $0$ for all $\fp$).

Since $K$ is torsion-free over $A$, the localization $M_\fp$ is identified with the $A_\fp$-submodule $A_\fp M=\{x/s\}$ of $K$: the relation $s''(s'm-sm')=0$ in $M\subseteq K$ with $s''\ne0$ is equivalent to $s'm=sm'$.

($\Leftarrow$) If $M_\fp$ is free over $A_\fp$ for every prime $\fp$, in particular for every maximal ideal $\fm$, then $M_\fm$ is flat over $A_\fm$ (Prop.~2.2(a)), so $M$ is flat over $A$ by Prop.~2.13 ((iii)$\Rightarrow$(i)).

($\Rightarrow$) Let $M\neq0$ be flat and $\fp$ a prime. Then $M_\fp$ is flat over $A_\fp$ (Prop.~2.13) and finitely generated, and $A_\fp$ is local, so $M_\fp$ is free by Thm.~2.16. Since $M\ne0$, pick $0\ne x\in M$; then $x/1\in M_\fp\subseteq K$ is nonzero, so $M_\fp\ne0$ and its rank is $\ge1$. If the rank were $\ge2$, there would be two elements $u,v$ of a basis of $M_\fp$, nonzero elements of $K$, say $u=a/b$ and $v=c/d$ with $a,b,c,d\in A\setminus\{0\}$. Then $(cb)u-(ad)v=ca-ac=0$ with $cb\neq0$, a nontrivial linear relation over $A_\fp$ (as $A_\fp\subseteq K$ is a domain), contradicting linear independence. So $M_\fp$ is free of rank $1$. \hfill$\square$

\exercise{2.10}{Let $A$ be an integral domain and $B$ its integral closure in the field of fractions $\operatorname{Frac}(A)$. Suppose $B$ is a finitely generated $A$-module. Show that $B$ is flat over $A$ if and only if $B=A$. (One can show that this is true without the finiteness assumption.)}
\pf
If $B=A$ then $B$ is free of rank $1$ over $A$, hence flat.

Conversely let $B$ be flat. Put $K=\operatorname{Frac}(A)$; then $A\subseteq B\subseteq K$, and $B\neq0$ is a finitely generated $A$-submodule of $K$. By Exercise~2.9, for every prime $\fp$ of $A$ the $A_\fp$-module $B_\fp=A_\fp B\subseteq K$ is free of rank $1$: $B_\fp=xA_\fp$ for some $x\in K^\times$. Since $1\in A\subseteq B_\fp$, we can write $1=ax$ with $a\in A_\fp$. Every element of $B_\fp$ is of the form $b/s$ with $b\in B$ integral over $A$ and $s\in A\setminus\fp$, hence is integral over $A_\fp$; in particular $x$ satisfies an equation $x^n+c_1x^{n-1}+\dots+c_n=0$ with $c_i\in A_\fp$ ($n\ge1$). Multiplying by $a^{n-1}$ and using $ax=1$:
\[ x+c_1+c_2a+\dots+c_na^{n-1}=0, \]
so $x\in A_\fp$. Then $x$ and $a=x^{-1}$ are both in $A_\fp$, so $x$ is a unit of $A_\fp$ and $B_\fp=xA_\fp=A_\fp$. Thus $(B/A)_\fp=B_\fp/A_\fp=0$ for every prime $\fp$ (localization is exact, Cor.~2.11), and in particular for every maximal ideal. By Lemma~2.12, $B/A=0$, i.e.\ $B=A$. \hfill$\square$

\exercise{2.11}{Let $M$ be a flat $A$-module, and let $\rho:A\to B$ be a ring homomorphism. Show that the following properties are true.
\begin{enumerate}[label=(\alph*),leftmargin=*]
\item If $A$ is an integral domain, then $M$ is torsion-free.
\item The canonical homomorphism $M\to M\otimes_A B$ is injective if and only if $\Ker(\rho)$ is contained in $\Ann(M)$.
\end{enumerate}}
\pf
(a) Let $a\in A$, $a\neq0$. Since $A$ is a domain, multiplication $t_a:A\to A$ is injective. As $M$ is flat, $t_a\otimes\operatorname{id}_M:A\otimes_A M\to A\otimes_A M$ is injective; under $A\otimes_A M=M$ this is multiplication by $a$ on $M$. So $ax=0$ implies $x=0$, i.e.\ $M$ is torsion-free.

(b) Put $I=\Ker(\rho)$. The sequence $0\to I\to A\to B$ is exact. Write it as $0\to I\to A\to \rho(A)\to 0$ followed by the injection $\rho(A)\to B$. Tensoring with the flat module $M$ gives an exact sequence $0\to I\otimes M\to M\to \rho(A)\otimes M$ and an injection $\rho(A)\otimes M\to B\otimes M$. Hence
\[ 0\to I\otimes_A M\to M\to M\otimes_A B \]
is exact, where the middle map is $x\mapsto x\otimes 1$. Therefore the kernel of $M\to M\otimes_AB$ is the image of $I\otimes_AM\to M$, which is $IM$. So the map is injective iff $IM=0$, iff $ax=0$ for all $a\in I$, $x\in M$, iff $I\subseteq\Ann(M)$. \hfill$\square$

\exercise{2.12}{Let $B$ be a flat algebra over a Dedekind domain $A$. Let $f\in B$ be such that for every maximal ideal $\mathfrak m$ of $A$, the image of $f$ in $B/\mathfrak mB$ is not a zero divisor. Show that $B/fB$ is still flat over $A$.}
\pf
By Cor.~2.14, a module over the Dedekind domain $A$ is flat iff it is torsion-free. In particular $B$ is torsion-free over $A$. We must show that $B/fB$ is torsion-free.

\textbf{Step 1: $f$ is a non-zero-divisor on $B/aB$ for every $a\in A$, $a\ne0$.} Let $K$ be the kernel of multiplication by $f$ on $B/aB$. Since localization is exact, $K_{\mathfrak m}$ is the kernel of $f$ on $B_{\mathfrak m}/aB_{\mathfrak m}$, and $K=0$ iff $K_{\mathfrak m}=0$ for all maximal $\mathfrak m$ (Lemma~2.12). Fix $\mathfrak m$ and put $A'=A_{\mathfrak m}$, $B'=B_{\mathfrak m}=B\otimes_AA'$. If $a$ is a unit of $A'$ (in particular if $A$ is a field), then $B'/aB'=0$ and there is nothing to prove. Otherwise $A'$ is a local principal ideal domain which is not a field, so its maximal ideal is $\pi A'$ with $\pi\neq0$ and every nonzero element is a unit times a power of $\pi$; thus $aB'=\pi^eB'$ for some $e\ge1$. Since $B'$ is torsion-free over $A'$ (flat), $\pi$ is a non-zero-divisor on $B'$. Moreover $B'/\pi B'=B\otimes_A A'/\mathfrak mA'=B/\mathfrak mB$ (elements of $A\setminus\mathfrak m$ already act invertibly on the $A/\mathfrak m$-module $B/\mathfrak mB$), and $f$ is a non-zero-divisor there by hypothesis.

We show by induction on $e\ge1$ that $f$ is a non-zero-divisor on $B'/\pi^eB'$; the case $e=1$ was just seen. Let $e\ge2$ and let $x\in B'$ with $fx\in\pi^eB'$. Then $fx\in\pi B'$, so $x\in\pi B'$ by the case $e=1$; write $x=\pi y$. Then $\pi fy=\pi^ez$ for some $z$, hence $fy=\pi^{e-1}z$ because $\pi$ is a non-zero-divisor. By induction $y\in\pi^{e-1}B'$, so $x\in\pi^eB'$. This proves Step 1.

\textbf{Step 2.} Let $a\in A\setminus\{0\}$ and $b\in B$ with $ab\in fB$, say $ab=fc$. Then $fc\in aB$, so $c\in aB$ by Step 1, say $c=ac'$. Thus $a(b-fc')=0$, and since $B$ is torsion-free, $b=fc'\in fB$. So the class of $b$ in $B/fB$ is zero: $B/fB$ is torsion-free, hence flat over $A$ by Cor.~2.14. \hfill$\square$

\exercise{2.13}{Let $M$ be a faithfully flat $A$-module. Let $N'\to N\to N''$ be a sequence of $A$-modules. Show that it is exact if and only if the sequence $N'\otimes M\to N\otimes M\to N''\otimes M$ is exact.}
\pf
Write the maps as $u:N'\to N$, $v:N\to N''$.

($\Rightarrow$) This is flatness: if $\Imm u=\Ker v$, then tensoring $0\to \Ker v\to N\to \Imm v\to0$ and $N'\to\Imm u\to 0$ with the flat module $M$ shows that $\Ker(v\otimes1)=\Ker v\otimes M=\Imm u\otimes M=\Imm(u\otimes 1)$ (all as submodules of $N\otimes M$, using that $\Ker v\otimes M\to N\otimes M$ is injective by flatness).

($\Leftarrow$) Assume the tensored sequence is exact. First, $vu=0$. Indeed let $h=vu:N'\to N''$, $h=\iota\circ p$ with $p:N'\to\Imm h$ onto and $\iota:\Imm h\to N''$ the inclusion. Then $h\otimes1=(\iota\otimes1)(p\otimes1)$ and $(v\otimes1)(u\otimes1)=0$ by assumption. Now $p\otimes1$ is surjective, and $\iota\otimes1$ is injective since $M$ is flat. Hence $\Imm h\otimes M=0$, and so $\Imm h=0$ by Prop.~2.17(ii).

So $\Imm u\subseteq\Ker v$; let $H=\Ker v/\Imm u$. By flatness, $\Ker v\otimes M$ and $\Imm u\otimes M$ are the submodules $\Ker(v\otimes1)$ and $\Imm(u\otimes1)$ of $N\otimes M$ (as in the first paragraph), and by right exactness $H\otimes M=(\Ker v\otimes M)/(\Imm u\otimes M)=\Ker(v\otimes1)/\Imm(u\otimes1)=0$. Hence $H=0$ by Prop.~2.17(ii), i.e.\ the original sequence is exact. \hfill$\square$

\exercise{2.14}{Let $A\to B$ be a ring homomorphism, and let $J$ be an ideal of $B$ such that $B/J$ is flat over $A$. Show that for any ideal $I$ of $A$, we have $(IB)\cap J=IJ$ (tensor the injection $I\to A$ by $B/J$).}
\pf
The inclusion $IJ\subseteq IB\cap J$ is clear. Conversely let $x\in IB\cap J$ and write $x=\sum_k i_kb_k$ with $i_k\in I$, $b_k\in B$. Put $\xi=\sum_k i_k\otimes b_k\in I\otimes_AB$ and let $\xi'=\sum_k i_k\otimes\bar b_k\in I\otimes_A(B/J)$ be its image. Since $B/J$ is flat, $I\otimes_A(B/J)\to A\otimes_A (B/J)=B/J$ is injective; it sends $\xi'$ to $\sum i_k\bar b_k=\bar x=0$ (as $x\in J$). Hence $\xi'=0$.

By right exactness of $I\otimes_A-$ applied to $J\to B\to B/J\to 0$, the sequence $I\otimes_AJ\to I\otimes_AB\to I\otimes_A(B/J)\to0$ is exact, so $\xi$ is in the image of $I\otimes_A J$: $\xi=\sum_l y_l\otimes c_l$ with $y_l\in I$, $c_l\in J$. Apply the multiplication map $I\otimes_AB\to B$, $i\otimes b\mapsto ib$: the image of $\xi$ is $\sum_k i_kb_k=x$ and also $\sum_l y_lc_l\in IJ$. Thus $x\in IJ$. \hfill$\square$

\exercise{2.15}{Let $0\to M'\to M\to M''\to0$ be an exact sequence of modules over a ring $A$. Let us suppose that $M''$ is flat. Show that $M'$ is flat if and only if $M$ is flat.}
\pf
Let $N_0\to N$ be an injective homomorphism of $A$-modules. Since $M''$ is flat, Prop.~2.6 gives commutative diagram with exact rows
\begin{align*}
0\to M'\otimes N_0&\to M\otimes N_0\to M''\otimes N_0\to 0\\
0\to M'\otimes N&\to M\otimes N\to M''\otimes N\to 0
\end{align*}
and vertical maps $\alpha',\alpha,\alpha''$ induced by $N_0\to N$ (so each arrow of the first row is followed down by the corresponding vertical map). The map $\alpha''$ is injective because $M''$ is flat.

($M$ flat $\Rightarrow$ $M'$ flat.) Let $x\in M'\otimes N_0$ with $\alpha'(x)=0$. Let $y$ be the image of $x$ in $M\otimes N_0$. Then $\alpha(y)$ is the image of $\alpha'(x)=0$, so $\alpha(y)=0$, and $\alpha$ is injective, so $y=0$. Since $M'\otimes N_0\to M\otimes N_0$ is injective (first row), $x=0$.

($M'$ flat $\Rightarrow$ $M$ flat.) Now $\alpha'$ and $\alpha''$ are injective. Let $y\in M\otimes N_0$ with $\alpha(y)=0$. The image of $y$ in $M''\otimes N_0$ is mapped by $\alpha''$ to the image of $\alpha(y)=0$, hence is $0$; so $y$ is the image of some $x\in M'\otimes N_0$. The image of $\alpha'(x)$ in $M\otimes N$ equals $\alpha(y)=0$, and $M'\otimes N\to M\otimes N$ is injective (second row), so $\alpha'(x)=0$, hence $x=0$ and $y=0$. So $\alpha$ is injective for every injection $N_0\to N$: $M$ is flat. \hfill$\square$

\exercise{2.16}{Let $(A,\mathfrak m)$ be a Noetherian local ring, and $C^\bullet: 0\to M'\to M\to M''\to0$ a complex of finitely generated flat $A$-modules. Show that if there exists an ideal $I\subseteq\mathfrak m$ such that $C^\bullet\otimes_AA/I$ is exact, then $C^\bullet$ is exact.}
\pf
By Thm.~2.16, the finitely generated flat modules over the local ring $A$ are free: $C^\bullet$ is $0\to A^a\xrightarrow{\varphi}A^b\xrightarrow{\psi}A^c\to0$ with $\psi\varphi=0$. We use repeatedly: if $N$ is a finitely generated $A$-module with $N=IN$, then $N=\mathfrak mN$ (as $I\subseteq\mathfrak m$), so $N=0$ by Nakayama (Lemma~2.7). Also $N\otimes_AA/I=N/IN$.

\textbf{$\psi$ is surjective.} $\Coker(\psi)\otimes A/I=\Coker(\psi\otimes A/I)=0$ since $C^\bullet\otimes A/I$ is exact. So $\Coker\psi=I\Coker\psi$, and $\Coker\psi=0$ by the remark above ($\Coker\psi$ is finitely generated). As $A^c$ is free, $\psi$ splits: $A^b\cong K\oplus A^c$ with $K=\Ker\psi$. In particular $K$ is a direct summand of $A^b$, hence finitely generated and projective, and the sequence $0\to K\to A^b\to A^c\to 0$ is split exact, so it stays exact after tensoring with $A/I$: $\Ker(\psi\otimes A/I)=K/IK$.

\textbf{$\varphi$ maps onto $K$.} Since $\psi\varphi=0$, $\varphi$ factors as $\varphi:A^a\to K$. By exactness of $C^\bullet\otimes A/I$, the induced map $\bar\varphi:(A/I)^a\to \Ker(\psi\otimes A/I)=K/IK$ is bijective. Then $\Coker(\varphi)\otimes A/I=\Coker\bar\varphi=0$, and $\Coker\varphi$ is finitely generated, so $\Coker\varphi=0$.

\textbf{$\varphi$ is injective.} As $K$ is projective and $\varphi:A^a\to K$ is onto, it splits: $A^a\cong \Ker\varphi\oplus K$. Hence $\Ker\varphi$ is finitely generated and $0\to\Ker\varphi\to A^a\to K\to0$ stays exact after $\otimes A/I$, so $\Ker\varphi\otimes A/I=\Ker\bar\varphi=0$. By the remark, $\Ker\varphi=0$.

Thus $\varphi$ is injective, $\Imm\varphi=K=\Ker\psi$, and $\psi$ is surjective: $C^\bullet$ is exact. \hfill$\square$

\exercise{2.17}{Let $M$ be a finitely generated flat $A$-module. Show that it is faithfully flat if and only if $\Ann(M)=0$.}
\pf
($\Rightarrow$) Let $J=\Ann(M)$ and $\iota:J\to A$ the inclusion. Since $M$ is flat, $\iota\otimes1:J\otimes_AM\to A\otimes_AM=M$ is injective; its image is $JM=0$. So $J\otimes_AM=0$, and then $J=0$ by Prop.~2.17(ii).

($\Leftarrow$) Suppose $\Ann(M)=0$; we verify Prop.~2.17(i). Let $\mathfrak m$ be a maximal ideal of $A$ and suppose $M=\mathfrak mM$. Localizing, $M_{\mathfrak m}=\mathfrak mM_{\mathfrak m}$ with $M_{\mathfrak m}$ finitely generated over the local ring $A_{\mathfrak m}$, so $M_{\mathfrak m}=0$ by Nakayama (Lemma~2.7). If $x_1,\dots,x_n$ generate $M$, each $x_i/1=0$ in $M_{\mathfrak m}$, so there are $s_i\in A\setminus\mathfrak m$ with $s_ix_i=0$. Then $s=s_1\cdots s_n\in A\setminus\mathfrak m$ kills $M$, i.e.\ $s\in\Ann(M)=0$. But $0\in\mathfrak m$, a contradiction. Hence $M\ne\mathfrak mM$ for all maximal $\mathfrak m$, and $M$ is faithfully flat. \hfill$\square$

\exercise{2.18}{Let $B$ be an $A$-algebra, and let $E$ be a faithfully flat $B$-module. Show that $E$ is flat over $A$ if and only if $B$ is flat over $A$.}
\pf
($\Leftarrow$) $E$ is flat over $B$ and $B$ is flat over $A$, so $E$ is flat over $A$ by Prop.~2.2(d).

($\Rightarrow$) Assume $E$ is flat over $A$. Let $N'\to N$ be an injective homomorphism of $A$-modules and let $g:N'\otimes_AB\to N\otimes_AB$ be the induced $B$-linear map; we must show $g$ injective. Using the natural isomorphisms $(N'\otimes_AB)\otimes_BE\cong N'\otimes_AE$ (and the same for $N$), the map $g\otimes_B\operatorname{id}_E$ becomes $N'\otimes_AE\to N\otimes_AE$, which is injective since $E$ is flat over $A$. Let $L=\Ker g$. Since $E$ is flat over $B$, $L\otimes_BE=\Ker(g\otimes_B\operatorname{id}_E)=0$. As $E$ is faithfully flat over $B$, $L=0$ by Prop.~2.17(ii). So $g$ is injective and $B$ is flat over $A$. \hfill$\square$

\exercise{2.19}{Let $f:A\to B$ be a faithfully flat ring homomorphism.
\begin{enumerate}[label=(\alph*),leftmargin=*]
\item Show that $f$ is injective and that $I\to I\otimes_AB$ is injective for every ideal $I$ of $A$.
\item Let $N=\Coker(f)$ be the cokernel of $f$. Let $I$ be an ideal of $A$. Using the commutative diagram
\begin{align*}
I&\to I\otimes_AB\to I\otimes_AN\to 0\\
0&\to I\to IB\to IN\to 0
\end{align*}
and Exercise 2.6(b), show that $I\otimes_AN\to IN$ is injective, and hence $N$ is a flat $A$-module.
\item Show that for any $A$-module $M$, the canonical map $M\to M\otimes_AB$ is injective.
\end{enumerate}}
\pf
(a) Let $K=\Ker f$. Since $B$ is flat, $K\otimes_AB\to A\otimes_AB=B$ is injective, and it sends $k\otimes b$ to $f(k)b=0$. Hence $K\otimes_AB=0$ and $K=0$ by Prop.~2.17(ii): $f$ is injective. Now let $I$ be an ideal of $A$. The composite $I\to I\otimes_AB\to B$, $x\mapsto x\otimes1\mapsto f(x)$, is injective since $f$ is; hence $I\to I\otimes_AB$ is injective.

(b) We first record Ex.~2.6(b): $IB\cap A=I$ (identifying $A$ with $f(A)\subseteq B$, allowed by (a)). Indeed let $P=(IB\cap A)/I$, the kernel of $A/I\to B/IB$. By flatness of $B$, $0\to P\otimes_AB\to (A/I)\otimes_AB\to (B/IB)\otimes_AB$ is exact. The composite $B/IB\cong (A/I)\otimes_AB\to(B/IB)\otimes_AB$ sends $\bar b$ to $\bar1\otimes b$; it has a left inverse $\bar c\otimes b\mapsto\overline{cb}$, so it is injective. Hence $P\otimes_AB=0$ and $P=0$, i.e.\ $IB\cap A=I$.

Identify $A$ with a subring of $B$, so $0\to A\to B\to N\to0$ is exact. Then $IN=(IB+A)/A\cong IB/(IB\cap A)=IB/I$, so the bottom row $0\to I\to IB\to IN\to0$ is exact. The top row $I\to I\otimes_AB\to I\otimes_AN\to0$ is exact by right exactness of $I\otimes_A-$. The vertical maps are the identity of $I$, the map $I\otimes_AB\to IB$, $x\otimes b\mapsto xb$, and the map $I\otimes_AN\to IN$, $x\otimes\bar b\mapsto x\bar b$; the diagram commutes, and the middle vertical map $I\otimes_AB\to IB$ is an isomorphism because $B$ is flat (it is the injection $I\otimes B\to B$ with image $IB$).

Let $z\in I\otimes_AN$ map to $0$ in $IN$. Lift $z$ to $w\in I\otimes_AB$. The image of $w$ in $IB$ maps to $0$ in $IN$, hence lies in (the image of) $I$: it is the image of some $x\in I$. Then $w$ and the image $w_x$ of $x$ in $I\otimes_AB$ have the same image in $IB$, so $w=w_x$ by injectivity of the middle map. Therefore $z$ is the image of $w_x$ in $I\otimes_AN$, which is $0$ since two consecutive maps of the top row compose to zero. Thus $I\otimes_AN\to IN$ is injective. Composed with $IN\subseteq N$ this says that $I\otimes_AN\to A\otimes_AN=N$ is injective for every ideal $I$, so $N$ is flat by Thm.~2.4.

(c) Let $M$ be an $A$-module. The sequence $0\to A\xrightarrow{f}B\to N\to0$ is exact and $N$ is flat by (b), so by Prop.~2.6 the sequence $0\to M\otimes_AA\to M\otimes_AB\to M\otimes_AN\to0$ is exact. In particular $M=M\otimes_AA\to M\otimes_AB$, $x\mapsto x\otimes1$, is injective. \hfill$\square$

\section*{1.3 Formal completion}

Conventions: rings are commutative with $1$; ``complete'' means separated and complete, as in the text (so $A$ is complete for the $I$-adic topology iff the canonical map $A\to\widehat{A}$ is bijective). For a topological group $G$ given by a filtration $(G_n)$, the closure of a subgroup $H$ is $\bigcap_n(H+G_n)$.

\exercise{3.1}{Is the usual topology on $\R$ defined by a subgroup filtration?}
\pf
No. Suppose $(G_n)_n$ is a descending chain of subgroups of $\R$ defining the usual topology, so the $G_n$ are neighborhoods of $0$ for the usual topology. Then $G_n\supseteq(-\varepsilon,\varepsilon)$ for some $\varepsilon>0$. For $x\in\R$ choose an integer $k\ge1$ with $|x/k|<\varepsilon$; then $x=k\cdot(x/k)\in G_n$ because $G_n$ is a group. Hence $G_n=\R$ for all $n$. But then the only neighborhood of $0$ for the topology defined by the filtration is $\R$ itself (the topology is indiscrete), whereas $(-1,1)$ is a neighborhood of $0$ for the usual topology. Contradiction. (The argument shows that no family of subgroups at all can give a fundamental system of neighborhoods of $0$ for the usual topology.) \hfill$\square$

\exercise{3.2}{Show that $\Z_p$ is a principal local ring, with maximal ideal $p\Z_p$.}
\pf
Recall $\Z_p=\varprojlim\Z/p^n\Z$ with $\widehat{G}_n=\{a=(a_m)\in\Z_p: a_m=0\text{ for }m\le n\}=\Ker(\Z_p\to\Z/p^n\Z)$.

\emph{Step 1: $\widehat{G}_n=p^n\Z_p$.} Clearly $p^n\Z_p\subseteq\widehat{G}_n$. Conversely let $a\in\widehat{G}_n$. For $m\ge n$, $a_m\in\Z/p^m\Z$ reduces to $0$ modulo $p^n$, so $a_m=p^nb_m$ for a unique $b_m\in\Z/p^{m-n}\Z$ (multiplication by $p^n$ is an injection $\Z/p^{m-n}\Z\to\Z/p^m\Z$ with image $p^n\Z/p^m\Z$). Uniqueness shows that $b_{m+1}\mapsto b_m$, so the family $b'_k=b_{k+n}\in\Z/p^k\Z$ ($k\ge1$) is an element $b'\in\Z_p$, and $p^nb'=a$ componentwise: the $k$-th component of $p^nb'$ is $p^nb_{k+n}$ reduced modulo $p^k$, which is the reduction of $a_{k+n}$ modulo $p^k$, namely $a_k$. Thus $\Z_p/p^n\Z_p\simeq\Z/p^n\Z$ and the topology of $\Z_p$ is the $p$-adic one; in particular $\bigcap_np^n\Z_p=0$ since $\bigcap_n\widehat{G}_n=0$.

\emph{Step 2: units.} If $u\in\Z_p\setminus p\Z_p$, then $u_1\ne0$ in $\Z/p\Z$, so each $u_m$ is a unit of $\Z/p^m\Z$ (it is nonzero modulo $p$). The inverses $u_m^{-1}$ are unique, hence compatible, and define an inverse of $u$ in $\Z_p$. Conversely elements of $p\Z_p$ are not units (their image in $\Z/p\Z$ is $0$ and $1\ne0$ there). So $\Z_p^\times=\Z_p\setminus p\Z_p$.

\emph{Step 3.} Let $a\ne0$. Since $a\in p^0\Z_p$ and $\bigcap_np^n\Z_p=0$, there is a largest $n$ with $a\in p^n\Z_p$; write $a=p^nu$, and $u\notin p\Z_p$ by maximality, so $u$ is a unit. Also $p^k\ne0$ in $\Z_p$ for every $k$ (its image in $\Z/p^{k+1}\Z$ is nonzero), hence $p^kw\ne0$ for each unit $w$. Therefore $\Z_p$ is an integral domain: $p^nu\cdot p^mv=p^{n+m}uv\ne0$.

Let $J\ne0$ be an ideal and let $n$ be the least integer such that $J$ contains an element $p^nu$ with $u$ a unit (every nonzero element has this form). Then $p^n\in J$, and every element of $J$ is $p^{k}w$ with $k\ge n$, so $J=p^n\Z_p$. Hence $\Z_p$ is a principal ideal domain. The non-units are exactly $p\Z_p$ (Step 2), an ideal, so $\Z_p$ is local with maximal ideal $p\Z_p$. \hfill$\square$

\exercise{3.3}{Let $A$ be the ring of germs of real analytic functions in $0\in\R$. Let $x\in A$ be the identity map on $\R$. Show that the map $f\mapsto\sum_n(f^{(n)}(0)/n!)T^n$ from $A$ to $\R[[T]]$ is injective, and induces an isomorphism from the $xA$-adic completion of $A$ onto $\R[[T]]$. Show that this is false if we replace $A$ by the $C^\infty$-functions in $0$.}
\pf
Write $\tau:A\to\R[[T]]$ for the Taylor map. It is a ring homomorphism (a germ of analytic function is, near $0$, the sum of its Taylor series, and the product of convergent series is the Cauchy product).

\emph{Injective.} If $\tau(f)=0$ then $f$ equals its Taylor series $0$ on a neighborhood of $0$, so the germ is $0$.

\emph{Powers of $xA$.} We claim $x^nA=\{f:\ \tau(f)\in T^n\R[[T]]\}$. The inclusion $\subseteq$ is clear. If $f=\sum_{k\ge n}a_kx^k$ converges on $(-r,r)$, then $g=\sum_{k\ge0}a_{k+n}x^k$ has the same radius of convergence, so it is a germ in $A$, and $f=x^ng$. Consequently $\tau$ induces an isomorphism $A/x^nA\to\R[[T]]/T^n\R[[T]]=\R[T]/(T^n)$: it is injective by the claim, and surjective since polynomials lie in $A$. These isomorphisms are compatible with the transition maps, so they induce an isomorphism
\[ \widehat{A}=\varprojlim A/x^nA\ \xrightarrow{\ \sim\ }\ \varprojlim\R[[T]]/T^n\R[[T]]=\R[[T]], \]
the last equality because a compatible family of classes mod $T^n$ is exactly a power series (Example 3.6 with $A=\R$, $r=1$). The composite $A\to\widehat{A}\to\R[[T]]$ is $\tau$, so the isomorphism is the one induced by $\tau$.

\emph{The $C^\infty$ case.} Let $E$ be the ring of germs at $0$ of $C^\infty$ functions. The Taylor map $E\to\R[[T]]$ is \emph{not} injective: $f_0(x)=e^{-1/x^2}$ ($f_0(0)=0$) is $C^\infty$, all its derivatives vanish at $0$, but it is not the zero germ. So the statement ``$f\mapsto\tau(f)$ is injective'' is false. Moreover, for $f\in E$ with vanishing derivatives of order $<n$, Taylor's formula with integral remainder $f(x)=x^n\int_0^1\frac{(1-t)^{n-1}}{(n-1)!}f^{(n)}(tx)\,dt$ shows that $f=x^ng$ with $g\in E$; hence again $E/x^nE\simeq\R[T]/(T^n)$ and the $xE$-adic completion is still $\R[[T]]$, but the canonical map $E\to\widehat{E}=\R[[T]]$ has kernel $\bigcap_nx^nE\ni f_0$ (it is surjective by Borel's lemma). So $E$ is not separated for the $xE$-adic topology, and $E\to\widehat{E}$ is not an isomorphism onto its image, in contrast with the analytic case. \hfill$\square$

\exercise{3.4}{Find examples where $M\otimes_A\widehat{A}\to\widehat{M}$ is not surjective, or is not injective.}
\pf
Take $A=\Z$, $I=p\Z$, so $\widehat{A}=\Z_p$ (Ex.~3.2). By Lemma 3.9 and Cor.~3.14 the map is surjective, resp.\ bijective, for finitely generated $M$, so $M$ must be infinitely generated.

\emph{Not surjective.} Let $M=\bigoplus_{k\ge1}\Z e_k$. Then $\widehat{M}=\varprojlim M/p^nM$ with $M/p^nM=\bigoplus_k(\Z/p^n)e_k$. Put $x_n=\sum_{k=1}^np^{k-1}e_k\in M/p^nM$. Since $x_{n+1}-x_n\equiv p^ne_{n+1}\equiv0\pmod{p^n}$, we get $x=(x_n)\in\widehat{M}$. The image of $M\otimes\Z_p=\bigoplus_k\Z_pe_k$ consists of elements $\sum_{k\le N}a_ke_k$ (finite sums, $a_k\in\Z_p$) whose $n$-th component has $e_{N+1}$-coordinate $0$ for all $n$. But the $e_{N+1}$-coordinate of $x_n$ is $p^N\ne0$ in $\Z/p^n$ for $n>N$. So $x$ is not in the image.

\emph{Not injective.} Let $M=\Q$. Then $p^n\Q=\Q$ for all $n$, so $\widehat{M}=\varprojlim\Q/\Q=0$. But $\Q\otimes_\Z\Z_p=S^{-1}\Z_p$ with $S=\Z\setminus\{0\}$; the map $\Z\to\Z_p$ is injective (a nonzero integer $m$ has nonzero image in $\Z/p^k\Z$ for $k$ large), so the image of $S$ in the domain $\Z_p$ (Ex.~3.2) does not contain $0$, and hence $\Z_p$ embeds in $S^{-1}\Z_p$; in particular $S^{-1}\Z_p\ne0$. So the map $\Q\otimes\Z_p\to\widehat{\Q}=0$ is not injective. \hfill$\square$

\exercise{3.5}{Let $B=\bigoplus_{n\ge0}B_n$ be a graded ring. (a) Show that $B_0$ is a subring of $B$, with unit. (b) Show that if $B$ is Noetherian, then so is $B_0$, and $B$ is a finitely generated algebra over $B_0$.}
\pf
(a) $B_0$ is an additive subgroup and $B_0B_0\subseteq B_0$. Write $1=\sum_de_d$ with $e_d\in B_d$. For $b\in B_n$, the equation $b=b\cdot1=\sum_dbe_d$ has $be_d\in B_{n+d}$; comparing components of degree $n$ in the direct sum gives $b=be_0$. Since $B$ is generated additively by homogeneous elements, $e_0$ acts as the identity on $B$, so $e_0=e_0\cdot1=1$. Thus $1\in B_0$ and $B_0$ is a subring with the same unit.

(b) Let $B_+=\bigoplus_{n\ge1}B_n$. It is an ideal (as $B_dB_e\subseteq B_{d+e}$) and $B=B_0\oplus B_+$, so the projection $B\to B_0$ is a surjective ring homomorphism with kernel $B_+$, i.e.\ $B_0\simeq B/B_+$. Hence $B_0$ is Noetherian when $B$ is.

The ideal $B_+$ is generated by the set of homogeneous elements of degree $\ge1$ (these span $B_+$ additively). Since $B$ is Noetherian, finitely many of them, say $x_1,\dots,x_r$ of degrees $d_i\ge1$, already generate $B_+$. Let $C=B_0[x_1,\dots,x_r]\subseteq B$. We show $B_n\subseteq C$ by induction on $n$. For $n=0$ it is clear. Let $n\ge1$ and $b\in B_n\subseteq B_+$; write $b=\sum_ic_ix_i$ with $c_i\in B$. Taking the degree-$n$ component, $b=\sum_i(c_i)_{n-d_i}x_i$, where $(c_i)_{n-d_i}\in B_{n-d_i}$ (interpreted as $0$ if $n<d_i$). Since $n-d_i<n$, these lie in $C$ by induction, so $b\in C$. Hence $B=C$ is a finitely generated $B_0$-algebra. \hfill$\square$

\exercise{3.6}{Let $A$ be an integral domain that is not a field, and let $K$ be its field of fractions. Show that the conclusion of Prop.~3.11 does not hold if we take $M=K$, $N=A$, and $I\ne0$.}
\pf
(The printed hypothesis ``$I\ne0$'' must be read as: $I$ is a nonzero \emph{proper} ideal; such ideals exist since $A$ is not a field. For $I=A$ the conclusion holds trivially, and for $I=0$ as well, so the restriction is necessary.)

Put $M_n=I^nK$. Choose $0\ne a\in I$. For any $y\in K$ we have $y=a^n\cdot(y/a^n)\in I^nK$, so $M_n=K$ for all $n$. The filtration $(M_n)$ is an $I$-filtration ($IM_n\subseteq M_{n+1}$) and it is stable: $M_{n+1}=K=IM_n$. Now $N=A$ and $M_n\cap N=A$ for all $n$. The filtration $(M_n\cap N)_n$ is stable iff $A=I\cdot A=I$ for $n\gg0$, which fails because $I\ne A$. So the conclusion of Prop.~3.11 fails. (The hypotheses ``$A$ Noetherian, $M$ finitely generated'' of Prop.~3.11 are both violated: $K$ is not finitely generated over $A$ when $A$ is Noetherian and not a field, and $A$ need not be Noetherian.) \hfill$\square$

\exercise{3.7}{Let $(A,\fm)$ be a Noetherian local ring. Show that $\bigcap_{n\ge0}\fm^n=0$. Give a counter-example with $A$ not Noetherian.}
\pf
By Cor.~3.13 (with $M=A$, $I=\fm$), $\bigcap_n\fm^n$ is the set of $x$ for which there is $\alpha\in\fm$ with $(1+\alpha)x=0$. In a local ring $1+\alpha$ is a unit for $\alpha\in\fm$ (otherwise $1+\alpha\in\fm$ and so $1\in\fm$). Hence $x=0$.

\emph{Counter-example.} Let $E$ be the ring of germs at $0$ of $C^\infty$ functions on $\R$ (Ex.~3.3). If $f(0)\ne0$ then $1/f$ is $C^\infty$ near $0$, so $f$ is a unit; hence $E$ is local with maximal ideal $\fm=\{f:f(0)=0\}$. By Taylor's formula in Ex.~3.3, $\fm=xE$ and $\fm^n=x^nE$. The germ $f_0=e^{-1/x^2}$ lies in every $\fm^n$ (indeed $f_0/x^n$, extended by $0$, is $C^\infty$) and is nonzero. So $\bigcap_n\fm^n\ne0$; $E$ is not Noetherian by the first part. \hfill$\square$

\exercise{3.8}{Let $A$ be a Noetherian ring, and $I,J$ ideals of $A$. Let $\widehat{A}$ be the $I$-adic completion of $A$ and $(A/J)^\wedge$ the completion of $A/J$ for the $(I+J)/J$-adic topology. Show that there is a canonical isomorphism $\widehat{A}/J\widehat{A}\simeq(A/J)^\wedge$.}
\pf
Let $M=A/J$, a finitely generated $A$-module. For $n\ge0$ we have $I^nM=(I^n+J)/J=((I+J)/J)^n$ (the $n$-th power of the image of $I$ in $A/J$ is the image of $I^n$, and $(I+J)/J$ is the image of $I$). So the $I$-adic topology of the $A$-module $M$ is the $(I+J)/J$-adic topology of the ring $A/J$, and the $I$-adic completion $\widehat{M}$ of the module $M$ is $(A/J)^\wedge$. By Cor.~3.14, the canonical map $M\otimes_A\widehat{A}\to\widehat{M}$ is an isomorphism. Finally
\[ M\otimes_A\widehat{A}=(A/J)\otimes_A\widehat{A}\simeq\widehat{A}/J\widehat{A}. \]
Composing gives the canonical isomorphism $\widehat{A}/J\widehat{A}\simeq(A/J)^\wedge$ (it sends the class of $\alpha=(\alpha_n)_n$ to the family of the images of $\alpha_n$ in $A/(I^n+J)$). \hfill$\square$

\exercise{3.9}{(\emph{$n$th root}) Let $n\ge2$ be an integer. Let $D=\Z[1/n]$. (a) Consider the polynomial $S=(1+T)^n-1\in D[T]$. Show that $D[[S]]=D[[T]]$ and that there exists an $f(S)\in SD[[S]]$ such that $1+S=(1+f(S))^n$. (b) Let $A$ be a complete ring for the $I$-adic topology, where $I$ is an ideal of $A$. Suppose that $n$ is invertible in $A$. Let $x\in I$. Show that there exists a unique continuous homomorphism $\varphi:D[[S]]\to A$ such that $\varphi(S)=x$. Conclude that there exists a $y\in I$ such that $1+x=(1+y)^n$. (c) Show, by giving an example, that statement (b) is false if $n$ is not invertible in $A$.}
\pf
\textbf{Lemma L.} Let $R$ be a ring and $g=c_1T+c_2T^2+\cdots\in R[[T]]$ with $c_1\in R^\times$. Then the substitution $\sigma:R[[X]]\to R[[T]]$, $\sum b_kX^k\mapsto\sum b_kg^k$ (a $T$-adically convergent sum, since $g^k\in T^kR[[T]]$) is a ring isomorphism.

\emph{Proof.} $\sigma$ is a ring homomorphism (it is the substitution homomorphism). We have $g^k=c_1^kT^k+(\text{higher terms})$. If $h=\sum b_kX^k\ne0$ with lowest nonzero coefficient $b_m$, then $\sigma(h)$ has $T^m$-coefficient $b_mc_1^m\ne0$ and no lower terms; so $\sigma$ is injective. For surjectivity let $H\in R[[T]]$. Construct $b_0,b_1,\dots$ inductively so that $H-\sum_{k<m}b_kg^k\in T^mR[[T]]$: this is true for $m=0$; if it holds for $m$, let $e$ be the $T^m$-coefficient of $H-\sum_{k<m}b_kg^k$ and put $b_m=ec_1^{-m}$; then $H-\sum_{k\le m}b_kg^k\in T^{m+1}R[[T]]$. Then $\sigma(\sum b_kX^k)-H\in\bigcap_mT^mR[[T]]=0$. \hfill(Lemma)

(a) Write $S=nT+\binom n2T^2+\dots+T^n$. Its coefficient of $T$ is $n$, a unit of $D$. By Lemma L (with $R=D$, $g=S$), $\sigma:D[[X]]\to D[[T]]$, $X\mapsto S$, is an isomorphism, which is the meaning of $D[[S]]=D[[T]]$. Let $f(X)=\sigma^{-1}(T)\in D[[X]]$. Its constant term equals the constant term of $\sigma(f)=T$, namely $0$, so $f\in XD[[X]]$, i.e.\ $f(S)\in SD[[S]]$. Moreover $\sigma((1+f)^n)=(1+T)^n=1+S=\sigma(1+X)$, and $\sigma$ is injective, so $(1+f(X))^n=1+X$ in $D[[X]]$.

(b) \emph{Existence.} Since $n$ is invertible in $A$, the ring map $\Z\to A$ extends uniquely to $D\to A$; we regard $D\to A$ as given. For $g=\sum b_kS^k\in D[[S]]$ put $\varphi_N(g)=\sum_{k<N}b_kx^k$. Since $x^k\in I^k$, $(\varphi_N(g))_N$ is a Cauchy sequence; as $A$ is complete (so separated) it has a unique limit $\varphi(g)$. Clearly $\varphi$ is additive and $\varphi(S)=x$. For $g=\sum b_kS^k$, $h=\sum c_kS^k$, $\varphi_N(gh)-\varphi_N(g)\varphi_N(h)=-\sum_{i,j<N,\ i+j\ge N}b_ic_jx^{i+j}\in I^N$; letting $N\to\infty$ gives $\varphi(gh)=\varphi(g)\varphi(h)$ (limits are unique and the ring operations are continuous). Also $\varphi(1)=1$. Finally $\varphi(S^Ng)=x^N\varphi(g)\in I^N$, so $\varphi(S^ND[[S]])\subseteq I^N$ and $\varphi$ is continuous.

\emph{Uniqueness.} Let $\psi$ be any ring homomorphism with $\psi(S)=x$ (continuity is not even needed). On $D$ it is the canonical map, since $\psi(n)=n$ is invertible, so $\psi(1/n)=n^{-1}$. Write $g=\sum_{k<N}b_kS^k+S^Ng_N$; then $\psi(g)-\varphi_N(g)=x^N\psi(g_N)\in I^N$. The same holds for $\varphi$, so $\psi(g)-\varphi(g)\in\bigcap_NI^N=0$ because $A$ is separated.

\emph{The root.} By (a), $f(S)=Sf_1$ with $f_1\in D[[S]]$. Put $y=\varphi(f(S))=x\varphi(f_1)\in I$. Applying $\varphi$ to $(1+f(S))^n=1+S$ gives $(1+y)^n=1+x$.

(c) Take $n=2$, $A=\Z/8\Z$, $I=2A$. Then $I^3=0$, so $A$ is complete and separated for the $I$-adic topology ($\widehat{A}=\varprojlim A/I^k=A$), but $2$ is not invertible in $A$. Take $x=2\in I$. If $3=1+x=(1+y)^2$ with $y\in I$, then $1+y$ is odd and its square is $\equiv1\pmod 8$ (odd squares are $1$ mod $8$), whereas $3\not\equiv1\pmod8$. So no such $y$ exists. \hfill$\square$

\exercise{3.10}{Let $k$ be a field of characteristic different from $2$, and $A=k[x,y]/(y^2-x^2(x+1))$. Let $\widehat{A}$ be the $\fm$-adic completion, where $\fm=(x,y)A$. Show that $\widehat{A}\simeq k[[u,v]]/(uv)$.}
\pf
Let $B=k[x,y]$, $\fn=(x,y)$, $F=y^2-x^2(x+1)$. By Example 3.6, $\widehat{B}=k[[x,y]]$ ($\fn$-adic). Apply Ex.~3.8 with $I=\fn$, $J=(F)$ ($B$ is Noetherian): the $(I+J)/J$-adic topology on $A=B/J$ is the $\fm$-adic one, so
\[ \widehat{A}\simeq k[[x,y]]/F\,k[[x,y]]. \]
\emph{Square root of $1+x$.} The ring $k[[x]]=\widehat{k[x]}$ is complete for the $(x)$-adic topology and $2$ is invertible in it (char $k\ne2$). By Ex.~3.9(b) (applied to $A=k[[x]]$, $I=(x)$, $n=2$, and the element $x\in I$) there is $w=1+y_0\in k[[x]]$, $y_0\in(x)$, with $w^2=1+x$. Then $w$ is a unit, $w=1+(\text{terms of degree}\ge1)$ and
\[ F=y^2-x^2w^2=(y-xw)(y+xw)\quad\text{in }k[[x,y]]. \]
\emph{Change of variables.} Put $u=y-xw$, $v=y+xw$, elements of $\fn k[[x,y]]$. We show that $k[[x,y]]$ is identified with $k[[u,v]]$ so that $u,v$ become the variables. Use the standard universal property (same proof as in Ex.~3.9(b)): for a ring $R$ complete for the $\mathfrak J$-adic topology and elements $r,s\in\mathfrak J$ there is a unique ring homomorphism $k[[U,V]]\to R$ with $U\mapsto r,V\mapsto s$ extending the given map from $k$.

Let $g=xw\in k[[x]]$; its coefficient of $x$ is $1$, so by Lemma L (Ex.~3.9), $\sigma_g:k[[X]]\to k[[x]]$, $X\mapsto g$, is an isomorphism; let $h=\sigma_g^{-1}(x)\in k[[X]]$, so $h(g(x))=x$, and also $g(h(X))=X$ (indeed $\sigma_g(g\circ h)=g\circ h\circ g=g\circ x=g=\sigma_g(X)$, and $\sigma_g$ is injective). Define
\[ \Phi:k[[U,V]]\to k[[x,y]],\ U\mapsto u,\ V\mapsto v;\qquad \Psi:k[[x,y]]\to k[[U,V]],\ x\mapsto h\!\left(\tfrac{V-U}2\right),\ y\mapsto\tfrac{U+V}2, \]
legitimate since $2$ is invertible and the images lie in the maximal ideal $(U,V)$, resp.\ $(x,y)$. Then, with $t=(V-U)/2$, $\Psi(u)=\frac{U+V}2-g(h(t))=\frac{U+V}2-t=U$ and likewise $\Psi(v)=\frac{U+V}2+t=V$. So $\Psi\Phi$ and the identity are homomorphisms $k[[U,V]]\to k[[U,V]]$ agreeing on $U,V$, hence equal by uniqueness. Conversely $\Phi\Psi(y)=\frac{u+v}2=y$, and $\Phi\Psi(x)=h\big(\frac{v-u}2\big)=h(xw)=h(g(x))=x$; so $\Phi\Psi=\mathrm{id}$ by uniqueness. Hence $\Phi$ is an isomorphism, with $\Phi(UV)=uv=F$. Therefore
\[ \widehat{A}\simeq k[[x,y]]/(F)\simeq k[[U,V]]/(UV).\]
\hfill$\square$

\exercise{3.11}{Let $A$ be a complete ring for the $I$-adic topology, where $I$ is an ideal of $A$. Let $(M_n)_{n\ge0}$ be $A$-modules such that $I^{n+1}M_n=0$ and that there exist surjective homomorphisms $\pi_n:M_{n+1}\to M_n$ with $\Ker(\pi_n)=I^{n+1}M_{n+1}$. Let $M=\varprojlim_nM_n$ and denote the (surjective) canonical homomorphisms by $u_n:M\to M_n$.
\begin{enumerate}[label=(\alph*),leftmargin=*]
\item Fix $d\ge0$. Show that for any $n\ge d$, there is an exact sequence $0\to I^{d+1}M_n\to M_n\to M_d\to0$.
\item Suppose $M_0$ is generated over $A$ by a finite number of elements $e_{0,1},\dots,e_{0,m}$. Let $e_1,\dots,e_m\in M$ be such that $u_0(e_i)=e_{0,i}$, and define $\varphi_n:A^m\to M_n$ by $(a_1,\dots,a_m)\mapsto\sum_ia_iu_n(e_i)$. Show that $\varphi_n$ is surjective, and, using Lemma 3.1, that $M$ is generated by the $e_i$.
\item Suppose moreover that $A$ is Noetherian and $M_0$ is finitely generated over $A$. Show that $\Ker u_n=I^{n+1}M$ (in other words, $u_n$ induces an isomorphism $M/I^{n+1}M\simeq M_n$).
In (c), the book proceeds in three steps: (1) Let $K_n=\Ker u_n$. Show that $K_n=I^{n+1}M+K_{n+1}$. (2) Show that the canonical homomorphism $\varprojlim_nM/I^{n+1}M\to\varprojlim_nM_n$ is an isomorphism (use Cor.~3.14). (3) Apply Lemma 3.1 to the exact sequence of inverse systems $0\to(K_n/I^{n+1}M)_n\to(M/I^{n+1}M)_n\to(M_n)_n\to0$ and show that $K_n=I^{n+1}M$.
\end{enumerate}}
\pf
(Parts (1)--(3) are written out as steps of (c).) Throughout, $\pi_{n,d}:M_n\to M_d$ ($n\ge d$) denotes the composite $\pi_d\circ\cdots\circ\pi_{n-1}$, so $\pi_d\circ u_{d+1}=u_d$ etc. The maps $u_n$ are surjective because the $\pi_n$ are (lift an element of $M_n$ successively to $M_{n+1},M_{n+2},\dots$).

(a) $\pi_{n,d}$ is surjective as a composite of surjections. We show by induction on $n\ge d$ that $\Ker\pi_{n,d}=I^{d+1}M_n$. For $n=d$, $\pi_{d,d}=\mathrm{id}$ and $I^{d+1}M_d=0$ by hypothesis. Suppose it holds for $n$. Then
\[ \Ker\pi_{n+1,d}=\pi_n^{-1}(\Ker\pi_{n,d})=\pi_n^{-1}(I^{d+1}M_n)=I^{d+1}M_{n+1}+\Ker\pi_n, \]
the last equality because $\pi_n$ is surjective, so $\pi_n(I^{d+1}M_{n+1})=I^{d+1}M_n$. Since $\Ker\pi_n=I^{n+1}M_{n+1}\subseteq I^{d+1}M_{n+1}$ for $n\ge d$, the right side is $I^{d+1}M_{n+1}$. This proves the exact sequence $0\to I^{d+1}M_n\to M_n\xrightarrow{\pi_{n,d}}M_d\to0$.

(b) \emph{Surjectivity of $\varphi_n$,} by induction on $n$. For $n=0$, $\varphi_0$ is onto since $u_0(e_i)=e_{0,i}$ generate $M_0$. Suppose $\varphi_n$ is onto and let $N=\varphi_{n+1}(A^m)\subseteq M_{n+1}$. Since $\pi_n\circ u_{n+1}=u_n$, $\pi_n(N)=\varphi_n(A^m)=M_n$, so $M_{n+1}=N+\Ker\pi_n=N+I^{n+1}M_{n+1}$. Substituting this into itself, $I^{n+1}M_{n+1}=I^{n+1}N+I^{2n+2}M_{n+1}$, and $I^{2n+2}M_{n+1}\subseteq I^{n+2}M_{n+1}=0$. Hence $M_{n+1}=N+I^{n+1}N=N$.

\emph{Generation of $M$.} Put $L_n=A^m/I^{n+1}A^m$. Since $I^{n+1}M_n=0$, $\varphi_n$ factors through a surjection $\bar\varphi_n:L_n\to M_n$, compatible with the transition maps ($L_{n+1}\to L_n$ the projection), so we have a surjection of inverse systems with kernels $K'_n=\Ker\bar\varphi_n$. The maps $K'_{n+1}\to K'_n$ are surjective: let $a\in A^m$ with class in $K'_n$. Then $\varphi_{n+1}(a)\in\Ker\pi_n=I^{n+1}M_{n+1}=I^{n+1}\varphi_{n+1}(A^m)=\varphi_{n+1}(I^{n+1}A^m)$, so there is $b\in I^{n+1}A^m$ with $\varphi_{n+1}(a-b)=0$; the class of $a-b$ lies in $K'_{n+1}$ and maps to the class of $a$ in $L_n$. By Lemma 3.1, $\varprojlim L_n\to\varprojlim M_n=M$ is surjective. Since $A$ is complete, $\varprojlim A^m/I^{n+1}A^m=\widehat{A}^m=A^m$ (the map $A^m\to\varprojlim L_n$ is bijective), and the composite $A^m\to M$ is $a\mapsto(\varphi_n(a))_n=\sum a_ie_i$. So $M=\sum Ae_i$.

(c) Now $A$ is Noetherian and $M_0$ finitely generated; by (b) $M$ is finitely generated.

(1) For $x\in M$, $u_n(I^{n+1}x)\subseteq I^{n+1}M_n=0$, so $I^{n+1}M\subseteq K_n$; and $K_{n+1}\subseteq K_n$ because $u_n=\pi_nu_{n+1}$. Hence $I^{n+1}M+K_{n+1}\subseteq K_n$. Conversely let $x\in K_n$. Then $u_{n+1}(x)\in\Ker\pi_n=I^{n+1}M_{n+1}=u_{n+1}(I^{n+1}M)$ (as $u_{n+1}$ is onto). So there is $z\in I^{n+1}M$ with $x-z\in K_{n+1}$, and $x\in I^{n+1}M+K_{n+1}$.

(2) Since $I^{n+1}M\subseteq K_n$, $u_n$ induces $M/I^{n+1}M\to M_n$, compatible with transitions, hence $\lambda:\varprojlim M/I^{n+1}M\to\varprojlim M_n=M$. Let $\iota:M\to\varprojlim M/I^{n+1}M=\widehat{M}$ be the canonical map. Then $\lambda\circ\iota=\mathrm{id}_M$. As $A$ is complete, $\widehat{A}=A$, so by Cor.~3.14 ($M$ is finitely generated over the Noetherian ring $A$) $\iota:M=M\otimes_AA\to\widehat{M}$ is an isomorphism. Hence $\lambda=\iota^{-1}$ is an isomorphism.

(3) The inverse systems $(K_n/I^{n+1}M)$, $(M/I^{n+1}M)$, $(M_n)$ with maps induced by the identity of $M$ (so $K_{n+1}/I^{n+2}M\to K_n/I^{n+1}M$ is induced by inclusion) give an exact sequence $0\to K_n/I^{n+1}M\to M/I^{n+1}M\to M_n\to0$. By (1), the image of $K_{n+1}$ in $K_n/I^{n+1}M$ is everything, so the transition maps of the first system are surjective. Lemma 3.1 gives an exact sequence
\[ 0\to\varprojlim K_n/I^{n+1}M\to\varprojlim M/I^{n+1}M\xrightarrow{\ \lambda\ }\varprojlim M_n. \]
Since $\lambda$ is injective by (2), $\varprojlim K_n/I^{n+1}M=0$. For an inverse system with surjective transition maps, each projection from the limit is onto (lift step by step), so every $K_n/I^{n+1}M=0$, i.e.\ $K_n=I^{n+1}M$. \hfill$\square$

\exercise{3.12}{Let $A$ be a Noetherian ring, $I$ an ideal of $A$, and $\widehat{A}$ the $I$-adic completion of $A$. Let $a\in A$. Show that if $a$ is not a zero divisor in $A$, then it is not a zero divisor in $\widehat{A}$. However, show, by giving an example, that $A$ can be an integral domain without $\widehat{A}$ being one.}
\pf
If $a$ is a non-zero-divisor, multiplication $a:A\to A$ is injective. Since $\widehat{A}$ is flat over $A$ (Thm.~3.15), tensoring with $\widehat{A}$ gives that multiplication by $a$ on $A\otimes_A\widehat{A}=\widehat{A}$ is injective, i.e.\ $a$ is not a zero divisor in $\widehat{A}$.

\emph{Example.} Let $k$ be a field of characteristic $\ne2$ (e.g.\ $\Q$), $A=k[x,y]/(F)$ with $F=y^2-x^2(x+1)$, and $I=(x,y)A$. $F$ is irreducible in $k[x,y]$: otherwise, being monic of degree $2$ in $y$ over the UFD $k[x]$, it would have a root in $k(x)$ (Gauss), i.e.\ $(y/x)^2=x+1$ in $k(x)$, impossible since $x+1$ has odd valuation at the prime $x+1$. So $(F)$ is prime and $A$ is a domain. By Ex.~3.10, $\widehat{A}\simeq k[[u,v]]/(uv)$, which is not a domain: $u\notin(uv)$, $v\notin(uv)$, but $uv=0$. \hfill$\square$

\exercise{3.13}{Let $A$ be a Noetherian ring, $I$ an ideal of $A$, and $M$ a finitely generated $A$-module. Suppose that $A$ is complete for the $I$-adic topology. Show that every submodule of $M$ is closed.}
\pf
Let $N\subseteq M$ be a submodule; we must show $\bigcap_n(N+I^nM)=N$. Put $M'=M/N$, a finitely generated $A$-module; the intersection is the preimage of $\bigcap_nI^nM'$, so it suffices to show $\bigcap_nI^nM'=0$.

Since $A$ is complete, $I$ lies in the Jacobson radical: for $\alpha\in I$ the series $\sum_{k\ge0}(-\alpha)^k$ converges in $A$ (partial sums are Cauchy) to an inverse of $1+\alpha$ (the product with $1+\alpha$ has partial sums $1-(-\alpha)^{N+1}\to1$, and $A$ is separated). By Cor.~3.13, an element $z\in\bigcap_nI^nM'$ satisfies $(1+\alpha)z=0$ for some $\alpha\in I$; as $1+\alpha$ is a unit, $z=0$. Hence $\bigcap_nI^nM'=0$ and $N$ is closed. \hfill$\square$

\exercise{3.14}{Let $A$ be a Noetherian ring, $I$ an ideal of $A$, and $\widehat{A}$ the formal $I$-adic completion of $A$. Show that the topology on $\widehat{A}$ induced by that on $A$ is the $I\widehat{A}$-adic topology.}
\pf
The topology on $\widehat{A}=\varprojlim A/I^n$ coming from that of $A$ is the one defined (Section 1.3.1) by the filtration $\widehat{A}_n=\Ker(p_n:\widehat{A}\to A/I^n)$. We show $\widehat{A}_n=I^n\widehat{A}$; since $(I\widehat{A})^n=I^n\widehat{A}$, the two filtrations coincide.

Tensor the exact sequence $0\to I^n\to A\to A/I^n\to0$ with $\widehat{A}$, which is flat (Thm.~3.15): the sequence $0\to I^n\otimes\widehat{A}\to\widehat{A}\to(A/I^n)\otimes_A\widehat{A}\to0$ is exact, and the image of $I^n\otimes\widehat{A}$ in $\widehat{A}$ is $I^n\widehat{A}$. So $I^n\widehat{A}=\Ker(\widehat{A}\to(A/I^n)\otimes_A\widehat{A},\ \alpha\mapsto1\otimes\alpha)$. Now $A/I^n$ is a finitely generated $A$-module whose $I$-adic completion is itself (it is discrete: $I^k(A/I^n)=0$ for $k\ge n$), and by Cor.~3.14 the canonical map $(A/I^n)\otimes_A\widehat{A}\to A/I^n$, $b\otimes\alpha\mapsto b\,p_n(\alpha)$, is an isomorphism. The composite $\widehat{A}\to(A/I^n)\otimes\widehat{A}\to A/I^n$ sends $\alpha\mapsto p_n(\alpha)$. Therefore $I^n\widehat{A}=\Ker p_n=\widehat{A}_n$. \hfill$\square$

\exercise{3.15}{(\emph{Mittag-Leffler condition}) Let $(A_n,\pi_n)_n$ be an inverse system of sets. For any pair $m>n$, let us denote the map $\pi_n\circ\pi_{n+1}\circ\cdots\circ\pi_{m-1}:A_m\to A_n$ by $\pi_{m,n}$. We say that the system $(A_n,\pi_n)_n$ satisfies the Mittag-Leffler condition if for every $n$, the descending sequence $(\pi_{m,n}(A_m))_{m>n}$ is stationary.
\begin{enumerate}[label=(\alph*),leftmargin=*]
\item Let $A'_n=\bigcap_{m>n}\pi_{m,n}(A_m)$. Show that $(A'_n,\pi_{n-1}|_{A'_n})_n$ is an inverse system and that the canonical map $\varprojlim_nA'_n\to\varprojlim_nA_n$ is bijective.
\item Suppose that $(A_n,\pi_n)_n$ satisfies the Mittag-Leffler condition and that $A_n\ne\emptyset$ for all $n$. Show that $A'_n\ne\emptyset$, that $A'_{n+1}\to A'_n$ is surjective, and that $\varprojlim_nA'_n\ne\emptyset$. Deduce from this that $\varprojlim_nA_n\ne\emptyset$.
\item Let $0\to(A_n)_n\to(B_n)_n\xrightarrow{\rho}(C_n)_n\to0$ be an exact sequence of inverse systems of Abelian groups such that $(A_n)_n$ satisfies the Mittag-Leffler condition. Let $(c_n)_n\in\varprojlim_nC_n$ and $X_n=\rho_n^{-1}(c_n)$, where $\rho_n:B_n\to C_n$ is the $n$-th component of $\rho$. Show that $(X_n)_n$ is an inverse system of sets satisfying the Mittag-Leffler condition. Deduce from this that $\varprojlim_nB_n\to\varprojlim_nC_n$ is surjective. This generalizes Lemma 3.1.
\end{enumerate}}
\pf
(a) Let $x\in A'_n$. Then $\pi_{n-1}(x)\in\pi_{n,n-1}(A_n)$, and for $m>n$ we can write $x=\pi_{m,n}(y)$ with $y\in A_m$, so $\pi_{n-1}(x)=\pi_{m,n-1}(y)\in\pi_{m,n-1}(A_m)$. Thus $\pi_{n-1}(x)\in A'_{n-1}$, and $(A'_n)$ is an inverse system. The canonical map $\varprojlim A'_n\to\varprojlim A_n$ is induced by the inclusions $A'_n\subseteq A_n$, hence injective. For surjectivity let $(a_n)\in\varprojlim A_n$. For all $m>n$, $a_n=\pi_{m,n}(a_m)\in\pi_{m,n}(A_m)$, so $a_n\in A'_n$ for every $n$, i.e.\ $(a_n)$ comes from $\varprojlim A'_n$.

(b) By hypothesis, for each $n$ there is $m(n)>n$ with $\pi_{m,n}(A_m)=A'_n$ for all $m\ge m(n)$ (the stationary value of the decreasing sequence is its intersection). As $A_m\ne\emptyset$, $A'_n\ne\emptyset$.

\emph{Surjectivity.} Let $x\in A'_n$. Take $m\ge\max(m(n),m(n+1))$. Then $x\in\pi_{m,n}(A_m)$, say $x=\pi_{m,n}(y)$, $y\in A_m$. Put $z=\pi_{m,n+1}(y)$; then $z\in\pi_{m,n+1}(A_m)=A'_{n+1}$ and $\pi_n(z)=\pi_{m,n}(y)=x$.

\emph{Nonempty limit.} Pick $a_0\in A'_0$; given $a_n\in A'_n$, surjectivity of $A'_{n+1}\to A'_n$ gives $a_{n+1}\in A'_{n+1}$ with $\pi_n(a_{n+1})=a_n$. Then $(a_n)\in\varprojlim A'_n\ne\emptyset$. By (a), $\varprojlim A_n\ne\emptyset$.

(c) Identify $A_n$ with a subgroup of $B_n$ (kernel of $\rho_n$). Write $\pi_n$ also for the transition maps of $B$, $C$; $\rho$ commutes with them and $\pi_n|_{A_{n+1}}$ is the transition map of $A$. If $b\in X_{n+1}$ then $\rho_n(\pi_n b)=\pi_n\rho_{n+1}(b)=\pi_n(c_{n+1})=c_n$, so $\pi_n(X_{n+1})\subseteq X_n$ and $(X_n)$ is an inverse system of sets. Each $X_n$ is nonempty ($\rho_n$ is onto) and is a coset of $A_n$: $X_n=x_n+A_n$ for any $x_n\in X_n$.

\emph{Mittag-Leffler.} For $m>n$, $\pi_{m,n}(X_m)=\pi_{m,n}(x_m)+\pi_{m,n}(A_m)$, a translate of the subgroup $\pi_{m,n}(A_m)\subseteq A_n$. Choose $m(n)$ such that $\pi_{m,n}(A_m)=A'_n$ for all $m\ge m(n)$. For $m\ge m(n)$ the sets $\pi_{m,n}(X_m)$ are cosets of the same subgroup $A'_n$ of $B_n$, and they form a decreasing sequence. If a coset $x+H$ is contained in a coset $y+H$ of the same subgroup, then $x\in y+H$ and $x+H=y+H$. So the sequence is constant for $m\ge m(n)$. Hence $(X_n)$ satisfies the Mittag-Leffler condition.

\emph{Surjectivity.} The $X_n$ are nonempty and satisfy Mittag-Leffler, so by (b) $\varprojlim X_n\ne\emptyset$. An element $(b_n)$ of it is a compatible family with $\rho_n(b_n)=c_n$, i.e.\ an element of $\varprojlim B_n$ mapping to $(c_n)$. Hence $\varprojlim B_n\to\varprojlim C_n$ is surjective. (If $A_{n+1}\to A_n$ is onto for all large $n$, then $(A_n)$ satisfies Mittag-Leffler, so this contains Lemma 3.1.) \hfill$\square$

\end{document}
